% This is the Reed College LaTeX thesis template. Most of the work 
% for the document class was done by Sam Noble (SN), as well as this
% template. Later comments etc. by Ben Salzberg (BTS). Additional
% restructuring and APA support by Jess Youngberg (JY).
% Your comments and suggestions are more than welcome; please email
% them to cus@reed.edu
%
% See http://web.reed.edu/cis/help/latex.html for help. There are a 
% great bunch of help pages there, with notes on
% getting started, bibtex, etc. Go there and read it if you're not
% already familiar with LaTeX.
%
% Any line that starts with a percent symbol is a comment. 
% They won't show up in the document, and are useful for notes 
% to yourself and explaining commands. 
% Commenting also removes a line from the document; 
% very handy for troubleshooting problems. -BTS

% As far as I know, this follows the requirements laid out in 
% the 2002-2003 Senior Handbook. Ask a librarian to check the 
% document before binding. -SN

%%
%% Preamble
%%
% \documentclass{<something>} must begin each LaTeX document
\documentclass[12pt,twoside]{reedthesis}
% Packages are extensions to the basic LaTeX functions. Whatever you
% want to typeset, there is probably a package out there for it.
% Chemistry (chemtex), screenplays, you name it.
% Check out CTAN to see: http://www.ctan.org/
%%
\usepackage{graphicx,latexsym} 
\usepackage{amssymb,amsthm,amsmath}
\usepackage{longtable,booktabs,setspace} 
\usepackage[hyphens]{url}
\usepackage{rotating}
\usepackage[numbers, sort]{natbib}
\usepackage{csquotes}
\usepackage{subfigure}
\usepackage[usenames,dvipsnames]{xcolor}
\definecolor{aapink}{RGB}{213,1,185}
\usepackage[colorlinks=true, linkcolor=aapink, urlcolor=aapink, citecolor=aapink]{hyperref}
\usepackage{float}
% Comment out the natbib line above and uncomment the following two lines to use the new 
% biblatex-chicago style, for Chicago A. Also make some changes at the end where the 
% bibliography is included. 
%\usepackage{biblatex-chicago}
%\bibliography{thesis}

% \usepackage{times} % other fonts are available like times, bookman, charter, palatino

\title{Icons and Age-Based Accessibility}
\author{Emily Shan-Yi Yu}
% The month and year that you submit your FINAL draft TO THE LIBRARY (May or December)
\date{May 2026}
\division{Mathematical and Natural Sciences}
\advisor{Advisor F. Name}
%If you have two advisors for some reason, you can use the following
%\altadvisor{Your Other Advisor}
%%% Remember to use the correct department!
\department{Computer Science}
% if you're writing a thesis in an interdisciplinary major,
% uncomment the line below and change the text as appropriate.
% check the Senior Handbook if unsure.
%\thedivisionof{The Established Interdisciplinary Committee for}
% if you want the approval page to say "Approved for the Committee",
% uncomment the next line
%\approvedforthe{Committee}

\setlength{\parskip}{0pt}
%%
%% End Preamble
%%
%% The fun begins:
\begin{document}

  \maketitle
  \frontmatter % this stuff will be roman-numbered
  \pagestyle{empty} % this removes page numbers from the frontmatter

% Acknowledgements (Acceptable American spelling) are optional
% So are Acknowledgments (proper English spelling)
    \hypertarget{ACKs}{\chapter*{Acknowledgments}}
    (a.k.a. ACKs, as CSCI389 has taught me)

    \paragraph{Thank you to: }
    \begin{itemize}
        \item \textbf{Dr. Adam Groce}, my thesis advisor and Computer Science professor: Thank you for a great year together (even if I spent a semester meeting you at 8am). Despite HCI not being one of your primary interests, I still appreciated your interdisciplinary insights a lot throughout this process! By the way, Ethics \& Policy is my favorite Reed course. You're the best!
        \item \textbf{Dr. Elizabeth Platte} of the Reed Instructional Technology Department: Thank you for such amazing insights into technological accessibility. Your higher-ed perspective truly helped inform the way I approached my own work, and I learned a lot about formatting this thesis accessibly (or as accessibly as possible) through your Document Accessibility course. Tell Susie I say thank you for a great past couple of years, what a coincidence!
        \item \textbf{Dr. Sabrina Schroerlucke and Dr. Greg Jensen} of the Reed Psychology Department: Thank you for your insights into experimental design and data analysis! I appreciate the time you dedicated to this non-Psychology thesising student.
        \item \textbf{Kayla Johnston, Claire Littman, and the rest of the IRB committee}: Thank you for your guidance throughout the IRB application process!
        \item \textbf{Dr. Michael Pearce} of the Reed Statistics Department: Thank you for your experimental statistics advice!
        \item \textbf{Dr. Josie Griffin and student DataLabby Leland Russell}: Thank you for your constant data analysis guidance! I had so many questions for your team and I appreciate the grace that you all had when answering me. 
        \item \textbf{Reed's Undergraduate Research Committee}: Thank you for financially supporting my research!
    \end{itemize}
    

% The preface is optional
% To remove it, comment it out or delete it.

    \chapter*{List of Abbreviations}
	\begin{table}[h]
	\centering % You could remove this to move table to the left
	\begin{tabular}{ll}
		\textbf{HCI}  	&  Human-Computer Interaction \\
		\textbf{AI} 	&  Artificial Intelligence\\
        \textbf{AIC}    & Akaike Information Criterion\\
        \textbf{ANOVA}  &  Analysis of Variance\\
        \textbf{ARIA}   &  Accessible Rich Internet Applications \\
        \textbf{DV}     & Dependent Variable\\
        \textbf{IO}     & Icon Only\\
        \textbf{IV}     & Independent Variable\\
        \textbf{IWTL}   & Icon with Text Labels\\
		\textbf{MTurk} 	&  Mechanical Turk\\
        \textbf{RT}     &  Reaction Time \\
        \textbf{TO}     & Text Only\\
        \textbf{UI}  	&  User Interface \\
        \textbf{USD} 	&  United States Dollars\\
        \textbf{UX}  	&  User Experience \\
		\textbf{WCAG}  	&  Web Content Accessibility Guidelines\\
		\textbf{WWW} 	&  World Wide Web\\
	\end{tabular}
	\end{table}
	

    \tableofcontents
% if you want a list of figures, also optional
    \listoffigures

% The abstract is not required if you're writing a creative thesis (but aren't they all?)
% If your abstract is longer than a page, there may be a formatting issue.
\chapter*{Abstract}
    In 2004, 70-98\% of websites were found to be inaccessible \cite{lazar_improving_2004}, and one of the user demographics most negatively affected by inaccessible websites is older users \cite{stephanidis_seven_2019}. Many human-computer interaction researchers have looked at whether different variants of Internet icons create positive performance effects for older users, with generally inconclusive results. 
    
    The goal of this thesis is to measure whether different icon types (specifically text, icons, or icons with text labels) have a positive impact on older users' performance (as compared to younger users) in a web-based task-completion setting. Specifically, users were assigned to an icon type and then shown six unique websites each containing one task. All research was conducted online and asynchronously. A total of 240 users completed the experiment. After processing data via a mixed-effects linear regression, the text-only condition performed the best, providing not only a boost in accuracy to both age groups, but specifically also a disproportionate boost in speed to the older user population. 



    \pagestyle{fancyplain} % turns page numbering back 

\chapter*{Introduction}
    In 2004, 70-98\% of websites were found to be inaccessible \cite{lazar_improving_2004}, and one of the user demographics most negatively affected by inaccessible websites is older users \cite{stephanidis_seven_2019}. One subject that researchers in the field of Human-Computer Interaction (HCI) study in order to create a more accessible Internet for older users is iconography. Even though semiotics has been a constant throughout human history (think languages), a unique lexicon of icons arose with the onset of the age of technology, defined heavily by Western culture and the newly emerging, globalization-influenced Internet culture \cite{ma_semiotics_1995}, making a brand new barrier to entry for older user populations.

    Many HCI researchers have looked at whether different variants of Internet icons create positive performance effects for older users, but results have generally been inconclusive. For instance, results varied on whether closer semantic distance, meaning “the closeness of meaning between icon and function” \cite{ganor_designing_nodate,liu_study_2023,dosso_how_2021,schroder_effects_2008}, or adding text labels \cite{wu_how_2022,dosso_how_2021,song_developing_2024} positively influences older user’s speeds and accuracies in ways significantly different than younger users. 

    Most of these studies were also conducted in cognitive psychology-inspired laboratory studies \cite{stephanidis_seven_2019}– they trade the web-based context the icons would be embedded into in the real world for isolated, causal inference. There are a couple of north star papers for this thesis that study icon accessibility within a more realistic context. Specifically, Yu \& Chattopadhyay presented 6 feature-rich mobile app user interfaces (UIs; e.g. Starbucks, AMC, Subway, etc.) and studied whether highlighting relevant features for the participants’s given task would boost performance speeds and accuracies among older adults\cite{yu_reducing_2024}. Excitingly, Yu \& Chattopadhyay found that when task-relevant features were emphasized, older adults performed at the same speeds and accuracies as younger adults during task completion \cite{yu_reducing_2024}. But accessibility iconography studies within task-based environments have been far less explored and do not have the wealth of findings that controlled lab experiment iconography studies have.
    
    Causal findings and realism are trade-offs. Yet, given the large strides in technological advancement being pushed, it’s now more important than ever to study accessibility within real-world applications. Alongside this, studying icons in the context of a realistic, task-based experimental design also could provide greater insight into the generally inconclusive findings of laboratory studies. 
    
 	The goal of this experiment is to measure whether different icon types (specifically text, icons, or icons with text labels) have a positive impact on older users' performance (as compared to younger users) in a web-based, task-completion setting. This experiment embedded these icons into six unique websites each containing one task, such as web browsing and food ordering. All research was conducted online and asynchronously. More specifically, participants were first recruited on Prolific, then redirected to a UX testing platform called Maze where the actual experiment was conducted, which tapped into Figma, a UI/UX design tool, to access these 6 webpages’ designs. Conducting experiments online a generally new procedure for many social sciences, and although HCI pioneers this push towards asynchronous research, accessibility research within HCI still primarily conducts studies in-person. With its novelty, a lot of unknowns exist and therefore a lot of scrutiny can be raised, which is addressed later in this paper in Sections \hyperlink{3}{3} and \hyperlink{4}{4}.
    
	A total of 240 users completed the experiment. Results were first examined as descriptive statistics and later plugged into a mixed-effects linear regression model to test for significance. The text-only condition came out significant, providing not only a boost in accuracy to both age groups, but also a disproportionate boost in speed to the older user population, which is the main finding of this paper. Alongside this, the icon with text labels condition overall performs worse than the text-only and icon-only conditions accuracy-wise.

    Although creating positive performance effects for both younger and older users is beneficial and should not be written off as null and void results, the exciting findings are the ones where older users are boosted is ways disproportionately greater than younger users. This is due to the idea of equity versus equality, where pushing for an equity-based design framework places older users on equal footing to allow them to access and experience technology in the same ways younger users do.





    
	
  \mainmatter % here the regular arabic numbering starts
  



















%The \introduction command is provided as a convenience.
%if you want special chapter formatting, you'll probably want to avoid using it altogether

\chapter{Background Literature}

In the following chapter, I will discuss the literature motivating my thesis. I will first begin in section \hyperlink{1.1}{1.1} by discussing the multidisciplinary origins of human-computer interaction (HCI) and how these origins inform HCI's current-day initiatives. Then section \hyperlink{1.2}{1.2} will focus in on the HCI subfield of age-based accessibility and the standards set for accessibility on the Internet. Section \hyperlink{1.3}{1.3} dives deeper into the academic research that motivates these standards, especially regarding computer iconography, and finally, section \hyperlink{1.4}{1.4} discusses task-based research involving older users in HCI and summarizes the gap in the literature that my thesis aims to address.

\hypertarget{1.1}{\section{Human-Computer Interaction}}

    Human-computer interaction (HCI) was born in the 1960s out of the minds of a handful of computer scientists during a time where computers were \enquote{scarce, expensive, bulky and formal-scheduled machines used for automatic calculations} \cite{roussel_looking_nodate}. Their goal was to make computers more compatible with humans to \enquote{\textit{augment the human intellect}} (emphasis in original), according to Douglas Engelbart, one of the founding fathers of HCI \cite{roussel_looking_nodate}. With this clear focus in mind, graphical user interfaces, mouses, applications, and much more were developed to make the computer human-usable \cite{myers_brief_1998}.

    In the 1980s, the tides began to turn as a wave of cognitive psychologists entered the field of HCI. With \enquote{[t]he body of knowledge, research findings and methods that made up cognitive psychology...} \cite{rogers_new_nodate}, HCI gained a robust toolkit and a broader, more holistic worldview.

    With the cognitive psychologists of the 1980s at the helm, the once narrowly-defined goal of HCI quickly expanded, leaving HCI \enquote{bursting at the seams} \cite{rogers_new_nodate}. Fast-forward to current day, HCI research includes everything from how individuals interact with their surrounding technology ecosystem (human-environment interactions), ethics in data privacy, alongside technology's potential in various sectors such as health, social governance, and education \cite{stephanidis_seven_2019}. 
    
    One of the biggest fields that HCI encompasses is accessibility, which \enquote{aims to ensure that IT applications and services can be used on an equal basis by users with disabilities and older users} \cite{stephanidis_seven_2019}. The need for expanded accessibility research cannot be overstated, as the U.S. Centers for Disease Control and Prevention estimated that one in four people in the United States live with a disability \cite{putnam_it_2023}. Accessibility is a subfield that, even within a very diverse parent field, stands out for its multidisciplinary basis. Beyond computer science and psychology, many other fields also inform HCI's accessibility research, including user interface/user experience (UI/UX) design \cite{putnam_it_2023}, ergonomics \cite{myers_brief_1998}, biology \cite{lindberg_how_2006, cruz_user-centered_2023}, economics \cite{nadkarni_task-based_2007}, and more.

\hypertarget{1.2}{\section{Older User Accessibility}}
    As Stephanidis et al.\@ explicitly mentions, one of the key subjects of HCI accessibility research is older users \cite{stephanidis_seven_2019}. After all, the National Institute of Aging states that people over 60 are one of the fastest growing user groups on the Internet \cite{kurniawan_research-derived_2005}, and older users' experience of the Internet vastly differs than that of younger users' \cite{stephanidis_seven_2019}. This is due to age-related declines in vision, which includes things like decreased color discrimination and narrower visual fields \cite{ganor_designing_nodate, cruz_user-centered_2023, lindberg_how_2006, noauthor_section508gov_nodate, noauthor_section508gov_nodate-1}, memory \cite{ganor_designing_nodate, noauthor_section508gov_nodate, noauthor_section508gov_nodate-1}, hearing \cite{noauthor_section508gov_nodate, noauthor_section508gov_nodate-1, cruz_user-centered_2023}, motility \cite{noauthor_section508gov_nodate, noauthor_section508gov_nodate-1, cruz_user-centered_2023, putnam_it_2023}, and more. Mahmud et al.~reveal that even the exploration strategies older users employ on technological interfaces are different than other age groups \cite{mahmud_learning_2020}. It should be noted that despite all of the biological developments of aging, older users and society as a whole might not see older individuals as \enquote{disabled}, which uniquely situates them in the wider accessibility discussion. For instance, some perceive this as a a natural phenomenon, whereas others find it hard to box these concurrent declines in ability levels into a singular disability.

    Yet despite this overwhelming need for accessibility, according to a 2004 paper, 70-98\% of websites are not accessible \cite{lazar_improving_2004}. In order to ease the burden of reaching in the dark to encode accessibility into websites, HCI researchers aim to update web accessibility standards to best match the needs of older users (and generally users with disabilities) today.
    
    The World Wide Web Consortium's Web Content Accessibility Guidelines (WCAG) are hailed as \enquote{the epitome of Web accessibility guidelines} \cite{kurniawan_research-derived_2005}. The WCAG is targeted towards software developers, allowing them to weave accessibility even at the granularity of individual lines of code. Recent amendments to the Americans with Disabilities Act also make the WCAG the definitive standard for any institution receiving federal funding, which extends to governmental organizations, higher education establishments, and more \cite{ada_fact_2024}. Other widely-used accessibility standards include Accessible Rich Internet Applications (ARIA), which is derived from UX standards \cite{putnam_it_2023}, and Section 508 of the U.S.'s Rehabilitation Act (29 U.S.C. 795d), which puts forth legal requirements for accessibility within American information and communication technology \cite{kurniawan_research-derived_2005, noauthor_section508gov_nodate, noauthor_section508gov_nodate-1, henry2006}. Other countries have also published their own guidelines on web accessibility \cite{kurniawan_research-derived_2005}. 

\hypertarget{1.3}{\section{Older Users and Icons}}
        But what informs these accessibility standards? Behind the scenes, HCI accessibility researchers have kept pace with the ever-changing realm of technology \cite{chong_bad_2023}, maintaining up-to-date standards even when innovation turns web accessibility on its head. One of these aspects of technology that has undergone substantial development in the mere decades since the inception of the Internet is its icons. 
        
        Iconography is a constant in human history. Semiotics, which is the study of icons, understands them as the foundation of human communication systems, with meanings that are socially and culturally imbued \cite{ma_semiotics_1995}. As the World Wide Web (WWW) emerged, its unique graphical interface naturally brought rise to the WWW's very own lexicon of iconography, defined primarily by Western culture and the newly emerging, globalization-influenced Internet culture \cite{myers_brief_1998, ma_semiotics_1995, knight_cultural_2009}. And as icons have emerged as a new mechanism of expressing meaning on the ever-changing Internet, HCI researchers have begun delving into the various facets of this brand new virtual iconography.

    \hypertarget{1.3.1}{\subsection{Semantic Distance}}
        One of these areas of interest is semantic distance, which is defined as \enquote{the closeness of meaning between icon and function} \cite{ganor_designing_nodate}. Results are varied on whether closer semantic distance has a significant positive influence on older users' speeds and accuracies. 

        Liu et al.~studied the impact of semantic distance (concrete versus abstract) and icon design (skeuomorphic versus flat, which is discussed in Section \hyperlink{1.3.2}{1.3.2}) on older versus younger participants' reaction times (RTs) \cite{liu_study_2023}. Specifically, Liu et al.~ended up studying four different icon types-- skeuomorphic concrete, flat concrete, skeuomorphic abstract, and flat abstract. They aimed to measure the resulting cognitive load of participants, which was quantified via performance measures such as reaction time, physiological measures (pupil diameter and the number of fixations), and subjective assessment measures. Participants were given semantic words, then were asked to judge whether the icon shown after the word matched. After analysis via a 3-factor analysis of variance (ANOVA), they found that when icons had further semantic distance \enquote{older participants had significantly higher reaction times... than younger participants} \cite{liu_study_2023}. 
        
        Dosso and Chevalier also found similar results in their study of semantic distance (close versus far) and text labels (icon versus label versus combination, which will be discussed in Section \hyperlink{1.3.3}{1.3.3}) on different age groups (young versus old) \cite{dosso_how_2021}. This study measured RTs (both including rereading and without rereading), correct response scores, and the number and length of rereads when presented with 30 interfaces of icons. Within the interface, there were 20 icons total, 19 being distractors and 1 being the target. Through ANOVAs performed per dependent variable, Dosso and Chevalier concluded that \enquote{older users needed more time to select icons, specifically when semantic distance was far} \cite{dosso_how_2021} as compared to younger users. 
        
        However, the sureness of these findings have been contested, such as by Ganor and Te'eni \cite{ganor_designing_nodate} and Schröder and Ziefle \cite{schroder_effects_2008}. Ganor and Te'eni observed whether differences in an icon's level of detail (simple versus detailed) and semantic distance (close versus distant) created any significant effects on different age groups' (young versus old) performance time and performance accuracy \cite{ganor_designing_nodate}. They achieved this by setting forth 40 tasks, where each task had participants choose the target icon amongst 7 other distractors. Although they did find a greater negative impact of semantic distance on performance on older participants, this impact was not significantly different than younger users according to their ANOVA analyses. 
        
        Schröder and Ziefle also looked at similar independent variables, manipulating icon concreteness and semantic distance (concrete versus abstract) and complexity (simple versus complex) on different age groups' (young versus old) RT and fraction of confirmatory responses per item (CR) \cite{schroder_effects_2008}. This resulted in 4 types of icons, namely abstract simple, abstract complex, concrete simple, and concrete complex icons. Schröder and Ziefle specifically showed mobile phone icons and gave 4 targets and 4 \enquote{dummies} per task for participants to select from. Through a 3-way repeated measures ANOVA, they found that although there was no significant effect of semantic distance on RT, there was a significant effect on CR, and still argued that concreteness is a more crucial factor than complexity when designing icons for elderly users despite no significant backing \cite{schroder_effects_2008}. 
        
        Generally, the consensus is that semantically close icons elicit faster reaction speeds, but this has not been confirmed to benefit elderly users more so than the average user in a significant manner.

    \hypertarget{1.3.2}{\subsection{Skeuomorphic Design}}

        As mentioned in Liu et al.\@, the design of icons have also been investigated with somewhat mixed results \cite{liu_study_2023}. Specifically, skeuomorphic icons are of great interest to HCI researchers. Skeuomorphic icons, or skeuomorphisms, are icons that \enquote{are usually designed with 3D effect graphics to imitate real objects through gradients, textures, highlights and shadows, and other visual features to bring users a realistic feeling to help them better understand the interface} \cite{liu_study_2023}. These are usually studied in conjunction with flat icons, both of which you can see in Figure \hyperlink{skeuo}{1.1}. Liu et al.\@, which as mentioned above had independent variables of semantic distance and skeuomorphic icon design \cite{liu_study_2023}, and Wu et al.\@, which studied skeuomorphisms in connection to text labels, among other conditions \cite{wu_how_2022}, both found that skeuomorphisms had a positive impact on older users' reaction speeds. 

        \begin{figure}[H]
            \hypertarget{skeuo}{}
            \centering
            \includegraphics[width=1\linewidth]{thesistemplate_2020-04-24/figures/skeuo.png}
            \caption{Examples of (a) Skeuomorphic Concrete, (b) Flat Concrete, (c) Flat Abstract, and (d) Skeuomorphic Abstract icons (reprinted by permission of Informa UK Limited, trading as Taylor \& Francis Group, \hyperlink{https://www.tandfonline.com}{https://www.tandfonline.com}) \cite{liu_study_2023}}
        \end{figure}
        
        Specifically, Wu et al.~looked at icon type (skeuomorphic icon and text, flat icon and text, skeuomorphic icon, and flat icon) and cognitive level on search accuracy, RT, and time to first fixation (the time it takes before glancing at the target icon) \cite{wu_how_2022}. Participants were all 60-66 years old and based in Hangzhou. They were asked to look at 3x3 arrays of icons, with 8 distractors and 1 target, all comprised of icons seen in medical services technology in the regions of Beijing, Shanghai, and Hangzhou. Skeuomorphisms elicited the highest search accuracies, fastest RTs, and fastest times to first fixation (with the order of icon types from best to worst in terms of search performance being skeuomorphic icons and text, flat icons and text, skeuomorphic icon, then flat icon) \cite{wu_how_2022}. These effects seemed to be the results of reduced semantic distance via realistic design. 

        But other studies, such as Xi and Wu, found that flat icons, as opposed to skeuomorphic icons, had higher search performance within older participant groups \cite{xi_influence_2018}. This was found in a study centered around icon style (line, metro, skeuomorphism, and flat) and its effect on total visit duration (per icon) and eye movement contrail (the moving line of sight during a task). Notably, this study's participant group were all graduate students, unlike all other studies within Sections \hyperlink{1.3}{1.3} and \hyperlink{1.4}{1.4}. Like Wu et al.\@, icons were arranged in a 3x3 grid, with the same ratio of distractors to target icons \cite{wu_how_2022}. After ANOVA analysis, they found that the total visit duration for flat icons was significantly shorter than skeuomorphic icons. Xi and Wu suggested could be due to the higher cognitive load which comes from deciphering the highly-detailed skeuomorphic designs. 

        Most results point towards skeuomorphic icons over flat icons, however Xi and Wu creates greater nuance in this debate with the introduction of results that show that flat iconography performed better than skeuomorphic icons for a participant pool of graduate students \cite{xi_influence_2018}.

\hypertarget{1.3.3}{\subsection{Icons and Text Labels}}

    A subtopic within accessible iconography that also has inconclusive findings is the impact of text labels. Previously mentioned studies such as Dosso and Chevalier \cite{dosso_how_2021} and Wu et al.~\cite{wu_how_2022}, among others, show that older users perform faster when presented icons with text labels as opposed to icons alone. However, Dosso and Chevalier state that this positive impact was experienced equally by both younger and older age groups \cite{dosso_how_2021}, which is seconded by Song et al.\@'s findings \cite{song_developing_2024}. 
    
    Song et al.\@'s study studied not only text labels, but also icon style (similar, symbolic, example, arbitrary), which mashes together the explorations of semantic distance (Section \hyperlink{1.3.1}{1.3.1}) and icon design (Section \hyperlink{1.3.2}{1.3.2}) \cite{song_developing_2024}. Participants of different ages (young versus old) were measured in RT and accuracy across 72 tasks, which, like Liu et al.\@, first showed the semantic name, then, like Wu et al.~and Xi and Wu \cite{wu_how_2022, xi_influence_2018}, showed a 3x3 array of 1 target and 8 distractor icons \cite{liu_study_2023}. According to a 2-way ANOVA, the RT of older participants increased across all icon types when text labels were provided, and most exciting of all, older and younger accuracies were equivalent after adding text labels to the icons (which without the text labels were not) \cite{song_developing_2024}. 

    Islam approaches this subtopic of iconography-based accessibility design in a slightly unconventional manner \cite{islam_exploring_2015}. Islam still conducts a laboratory study, but instead of providing analysis just via quantitative measures, Islam centers their behavioral study around semi-structured interviews. Specifically, after sourcing icons from 2 web domain interfaces, namely a museum webpage and a university webpage, they show the icons to the participants and ask them a series of both open-ended, probing, and closed questions, resulting in interviews between 100 to 120 minutes long. These interviews included questions such as: \enquote{What could be the referential meaning [of the icon]?} and \enquote{How confident are you that you are correct?} \cite{islam_exploring_2015}. Islam's participants were largely graduate students, but not entirely graduate students like \cite{xi_influence_2018}, with 22 out of 28 being graduate students, 5 out of 28 being full-time company workers, and 1 out of 28 being a research assistant. Through qualitative and quantitative analyses, Islam found that \enquote{the textual interface signs showed better performance in 4 groups; iconic signs showed better performance in 3 groups; and icon[s] with text sign[s] showed better performance in [the] remaining 4 groups} \cite{islam_exploring_2015}, pointing towards no definite right choice for whether text labels on icons are more accessible across the board for all age groups. 

    So all in all, text labels' impact has been shown to be beneficial for older users, but perhaps not in ways that disproportionately benefit the age group. Islam further asserts this by that icons with text labels do not have an universally positive performance impact. However, this icon-and-text-label interaction has not been widely studied, so this research aims to further add to this budding discussion on the impacts of text labels on icon accessibility for older users.

\hypertarget{1.4}{\section{Older Users and Task-Based Experimental Design}}
    The vast majority of the studies cited in section \hyperlink{1.3}{1.3} are modeled after traditional cognitive psychology methodology, recording reaction times to stimuli presented in isolation without the context of an actual website surrounding said stimuli \cite{stephanidis_seven_2019}. Isolating the icons does reap strong causal conclusions, but in the vast, virtual universe of the Internet, are icons ever shown in isolation? Causal findings and realism are trade-offs, and given the previously inconclusive laboratory experiments on older users and iconography, studying icons in the context of a realistic, task-based experimental design could provide greater additional insight into this ongoing conversation.

    And research on the negative impacts of website clutter on older user performance speeds further this point. Nadkarni phrases it best: \enquote{perceived website complexity (PWC) is central to understanding how sophisticated features of a website... affect a visitor's experience} \cite{nadkarni_task-based_2007}. Grahame et al.~defines clutter as pixels as a percentage of the available page space that was occupied by text or graphics \cite{grahame_age_2004}. This was primarily studied by controlling for the number of links on a webpage, but Grahame et al.~also manipulated the font size of the links, the placement of the target link on the UI, etc.. The two age groups (young versus old) were shown web pages modeled off of free-to-use templates and were recorded for number of errors, RT, and fixation number. Through ANOVAs calculated per dependent variable, Grahame et al.~determined that increased clutter significantly increased search times and increased error rates for older adults \cite{grahame_age_2004}. 

    And although HCI primarily conducts research via laboratory studies, there is a pocket of HCI that values task-based procedural research. One of the earliest examples of this can be found in 1993, where with the motivation of studying tasks women were \enquote{more likely to perform... in the real world} \cite{czaja_age_1993}, Czaja and Sharit studied how age impacted performance speeds for women, ages 25-75 years old, across 3 different tasks-- namely, data entry, file modification, and inventory management. Data entry was simple and repetitive, requiring the participants to change addresses according to instruction cards. File modification was an intermediary level task, where participants were given a database of high-schoolers organized by grade level and asked to check for consistency between individual student records and instruction cards. Finally, inventory management was the least structured and most demanding task, requiring participants to track the inventory of products from outside suppliers and manipulate this relational database according to instruction cards. Czaja and Sharit showed that increased age, as well as a lack of computer experience, were the greatest determiners of slower response times and a larger number of errors across all 3 tasks \cite{czaja_age_1993}. 
    
    In Mahmud et al.\@, three age groups (children, adults, and older adults) were studied, each with 10 participants per group \cite{mahmud_learning_2020}. However, this study takes a different approach than Czaja and Sharit-- Mahmud et al.~instead observed the exploratory styles of these 3 age groups, who all were first-time users of Microsoft's OneNote, when asked to complete 4 main tasks using the application \cite{mahmud_learning_2020, czaja_age_1993}. These tasks were complex, multistep tasks, including making a new notebook, adding relevant notes to a notebook, and sharing notebooks with others, etc.. After coding video footage of the participants using OneNote and analyzing pre- and post-exploration interviews, Mahmud et al.~found a great treasure trove of findings on older users in embedded, task-based exploration settings. For instance, as many of the other readings mentioned may suggest \cite{liu_study_2023, dosso_how_2021, grahame_age_2004}, older participants were overall slower at task completion as compared to other age groups. Their exploration style tended to not fully explore all options, rely on previous experience with other UIs to their own detriment, and older users experienced more breakdowns when faced with errors. In fact, older adults typically spent a longer time repeating failed sequences of selections, whereas adults and children more often redirected their efforts based on error messaging. And generally, Mahmud et al.~found that older participants were generally dissatisfied with self-exploration, whereas both children and adults maintained high satisfaction levels and confidence before and after the study. 

    The study that inspired my thesis the most is Yu and Chattopadhyay's 2024 study of mobile tasks within the context of complex, feature-rich UIs \cite{yu_reducing_2024}. This experiment involved multiple phases, first beginning with a young versus old (each with 30 participants) comparison of task completion given 4 relevant feature highlighting techniques, where weighted zoom and highlight with context prevailed in terms of task completion time. Yu and Chattopadyay found that with emphasized task-relevant features, older adults performed at the same speeds and accuracies during task completion as younger adults \cite{yu_reducing_2024}. And it wasn't that the older adults were not initially erroring as much when encountering a difficult task, it was that with the help of search space reduction, these older adults were recovering successfully, as opposed to the behaviors exhibited by older adults when they encountered errors in Mahmud et al.~\cite{mahmud_learning_2020}. The final phase was implementing an AI-assistant (Nav Nudge) to reduce the search space, which exhibited 94\% accuracy for their \textit{Query} dataset and 93\% accuracy for their \textit{UILabel} dataset. Yu and Chattopadyay presented the UI of a map application none of the participants (who were all older (10 participants)), none overlapping with the first phase of the study) were familiar with and studied whether providing Nav Nudge, which could decipher speech prompts and highlight the top three most relevant features based on the participant's given task would boost performance speeds and accuracies among older adults (via weighted zoom or highlight with context, as determined by the first part of this study). For example, one of the tasks assigned was: \enquote{Find out how long it will take to travel from Chicago O’Hare Airport to Chicago Union Station using public transit} \cite{yu_reducing_2024}. Specifically, these older adults were first presented with the UI without visual cues, and if they remained stuck on the given task for more than 13 seconds or willingly gave up, they were able to use Nav Nudge to attempt the task once more. As this phase was more exploratory, a qualitative analysis was performed to reveal that Nav Nudge successfully enabled older users in the task completion setting and that older users rated Nav Nudge highly \cite{yu_reducing_2024}. 
    
    If this decomplexification effort can be replicated with icons, whose explicit purpose is to use condensed symbols to represent larger meanings \cite{ma_semiotics_1995}, especially when studied in a task-based experimental design, this would be monumental for older users. As seen in the HCI research above, it is exceedingly rare to find design changes that significantly allow older users and younger users to perform as equals. If substantiated, these findings could introduce a brand new paradigm of Internet iconography and older user UX accessibility.

        
\chapter{Methodology}

This chapter discusses the methodology of this study. As an overview, all research was conducted online and asynchronously. This was achieved via a participant matching platform called Prolific for recruitment that then redirected participants to Maze, a UX usability testing platform, to measure speeds and accuracies in completing six tasks embedded within web interfaces. 

Section \hyperlink{2.1}{2.1} first justifies the use of Prolific for participant recruitment, especially in light of recent developments in artificial intelligence (AI) and its potential harm towards sound online experiment hosting. Then, Section \hyperlink{2.2}{2.2} runs through the experimental design (Prolific and Maze) and discusses the design choices behind the Figma UI wireframes used. Finally, Section \hyperlink{2.3}{2.3} discusses the data analysis strategy for both descriptive statistics and linear regression calculations.

\hypertarget{2.1}{\section{Prolific and Online Survey Platforms}}
    Prolific, as the introduction to this chapter states, is an online participant recruitment platform. It allows researchers to post studies to the platform for its exceptionally large user base of over 200,000 active users to participate in \cite{run_prolific}. Surveys can gather large samples in a matter of days or even hours. However, the reliability of the platform has seen significant doubt, especially as AI makes rapid gains in its autonomous decision-making and action-taking abilities.

\hypertarget{2.1.2}{\subsection{In Favor of Online Survey Platforms}}
    
    I begin by introducing a quote from an analysis of Prolific, which at the time was an up-and-coming participant recruitment platform: 
        \enquote{The success of online experiments is not surprising, as they offer at-scale recruitment of participants in a short time, are generally cheap, and offer access to a broader population – potentially even representative of the internet population – than classical lab experiments.} \cite{palan_prolificacsubject_2018}
        
    And these benefits are especially exciting in the age-based accessibility area of HCI, where recruiting sizable populations of older participants is particularly difficult. For instance, Yu and Chattopadhyay performed a power analysis for the first leg of their study which informed their total of 60 participants \cite{yu_reducing_2024}, and Mahmud et al.\@ had only 30 total participants \cite{mahmud_learning_2020} -- and as a reminder, these participant pools are further divided by age group, with Yu and Chattopadhyay having older and younger age groups and Mahmud et al.\@ having a children, adults, and older adults three-way split. And although the results in these papers are exciting, with such limited sample sizes, it's hard to ensure the generalizability of these findings and the representative diversity of these samples. This is mitigated with Prolific, which has a total \enquote{active, verified, and engaged} participant base of over 200,000 users \cite{run_prolific}.
    
    The cost-effectiveness of Prolific is also a key benefit, but it's of note that Prolific pays an ethical wage, unlike other platforms such as Amazon's Mechanical Turk (MTurk), which is known for \enquote{unreasonably low subject payment} \cite{palan_prolificacsubject_2018}. Specifically, Prolific ensures a minimum fixed payment per unit of time, calculated based on a minimum hourly reward of \$8 U.S. Dollars (USD), with a recommended hourly reward of \$12 USD \cite{how_prolific}. 

% might not actually wanna add... still kind of low lol
%And even though the sample size for this study goes significantly beyond what has previously been done in this field (200 total, 100 per age group), compensating each individual fairly does not break the bank-- even with the \$12 USD recommended rate, given the short estimated length of this experiment (5 minutes), each person gets \$1 USD for participating, which with Prolific fees ended up costing \$266.40 USD.

    Online survey platforms such as Prolific are no silver bullet, though. There's longstanding belief that these participants, in order to make cash out quick, speed through experiments to earn their incentive and move on to the next survey \cite{palan_prolificacsubject_2018}. This fear is only amplified with the onset of AI in an online survey context. Peterson and Asher et al.\@ both studied AI-generated free responses on MTurk and Prolific respectively \cite{peterson_impact_nodate, asher_chatbots_2025}. Peterson found that just 20\% of the 1443 responses passed quality checks and counted as usable data \cite{peterson_impact_nodate}, whereas Asher et al.'s findings on Prolific was a little more hopeful-- when using a keystroke tool to flag submissions with an unusually low number of keystrokes, they found that only 9\% were AI-assisted responses \cite{asher_chatbots_2025}. But generally, MTurk is more vulnerable to AI responses \cite{christoforou_generative_2024}, partially due to its unethical wages, which encourages long-tail exploitation that oftentimes isn't valuable enough to humans to engage in \cite{carlini_llms_2025}. And now, with the rise of commercially-available AI models that can easily execute computer tasks, such as open-source OpenClaw \cite{openclaw}, OpenAI's Operator \cite{openai_operator}, Claude's Sonnet 3.0 \cite{claude_sonnet}, and Google's Gemini 2.5 \cite{google_gemini}, things look even more bleak for online survey models such as Prolific. Given their novelty, little research has been done on the topic of so-called \enquote{computer use} AI models participating in online surveys yet. It's worth noting that, however, these autonomous agents run locally, meaning there is a slightly higher barrier to entry as compared to web-hosted AI chatbots.

    Luckily, Prolific takes measures against AI and participants who rush through studies. Prolific allows researchers to manually approve study participants for payment, which can help researchers screen for participants with unusually fast completion times or other suspicious activity \cite{prolific_home}. During the writing of this paper, Prolific is beta-ing an option to automatically \enquote{[r]eject exceptionally fast submissions} and to detect AI-generated written responses, which also deters against these bad actors \cite{prolific_home}. But none of these absolutely destroy the threat. In contemplating trade-offs, both via reading the literature presented above and discussing with Reed faculty and staff (such as Psychology Professor Greg Jensen, who has published research employing Prolific for recruitment \cite{jensen_does_2022}), the ability to seamlessly recruit a notably difficult demographic at such large numbers outweighed the threat of falsified data. If anything, by bulking up the sample size, it mitigates the effect of a handful of AI-submitted-responses, if any. But this is not to say that the risk is zero, which is a hard truth kept front-of-mind throughout the following phases of this thesis. 

    Another concern lies in whether the results will be overall faster than average because some online survey users tend to rush through surveys to earn more, and that generally these users will be more fluent in computer technology given that they actively participate on online survey platforms such as Prolific. This concern is valid, however, this is a bias that could have just as easily occurred in an in-person study as well, althought perhaps not to the same extent. To gather a significant number of older users, one could suggest recruiting them in community spaces elderly folks frequent, such as the library and senior living communities, but people who would willingly participant may be more confident and skilled in computer-use than the average older user. And in terms of diversity and quantity, online recruitment wins over local recruitment. 

    All in all, Prolific is not the perfect recruitment tool. But in order to get significant results and to ensure smooth recruitment and data collection, this study has decided to employ Prolific to meet its ends.

\hypertarget{2.2}{\section{Experimental Design}}
\hypertarget{2.2.1}{\subsection{Prolific Recruitment and Screening}}
    Participants were recruited via Prolific, the online survey platform discussed in Section \hyperlink{2.1}{2.1}. Prolific sought out participants who: 
    \begin{enumerate}
        \item Had 1+ years of web-use experience
        \item Had conversational English fluency
        \item Were located within the U.S.A. while completing this experiment
        \item Used a laptop or desktop device while completing this experiment
        \item Were an age that fell within either the younger (18-35 years old) or older (60+ years old) age groups
    \end{enumerate}
    
    As a sidebar, within HCI, \enquote{old} is fluid and doesn't refer to a single, well-defined age group. Some researchers define \enquote{old} by their nation's retirement ages, whereas others define old age via years in the workforce, and most rely on local community's perceptions to inform their definitions \cite{yu_reducing_2024}. These age groups are inspired by the age groups presented in the previous literature mentioned in Chapter \hyperlink{1}{1}, which had no clear consensus on what \enquote{young} and \enquote{old} mean, as seen in Figure \hyperlink{agetab}{2.1}. 

        \begin{figure}[H]
            \hypertarget{agetab}{}
            \centering
            \begin{tabular}{| p{3.5cm} | p{3.5cm} | p{3.5cm} | p{3.5cm} |}
                \hline
                 Study & Young Age Range (Years) & Old Age Range (Years) & Notes\\
                 \hline
                 Liu et al.~\cite{liu_study_2023} & 22-25 & 58-69 & \\
                 \hline
                 Dosso and Chevalier \cite{dosso_how_2021} & 18-24 & 60-81 & \\
                 \hline
                 Ganor and Te'eni \cite{ganor_designing_nodate} & 20-35 & 65+ &\\
                 \hline Schröder and Ziefle \cite{schroder_effects_2008} & 19-29 & 50-65 & \\
                 \hline
                 Wu et al.~\cite{wu_how_2022} & & 60-66 & Only studied older users\\
                 \hline
                 Xi and Wu \cite{xi_influence_2018} & & & Only studied graduate students\\
                 \hline 
                 Song et al.~\cite{song_developing_2024} & Mean age 22.33 & Mean age 67.92 & \\
                 \hline
                 Islam \cite{islam_exploring_2015} & 22-41 & & Did not study older users\\
                 \hline
                 Grahame et al.~\cite{grahame_age_2004} & 20-31 & 55-62 &\\
                 \hline
                 Mahmud et al.~\cite{mahmud_learning_2020} & 20-35 & 60-75 & Also studied children ages 10-12\\
                 \hline
                 Yu and Chattopadhyay \cite{yu_reducing_2024} & 20-29 & 60+ &\\
                 \hline
            \end{tabular}
            \caption{Table of age ranges in Section \hyperlink{1.3}{1.3} and \hyperlink{1.4}{1.4}'s readings}
            \label{fig:placeholder}
        \end{figure}
        
    And despite the concerns discussed in Section \hyperlink{2.1}{2.1}, it still felt necessary to confirm that they were proficient in computer technology, especially since this study centers older users, who have various impediments that negatively impact their ability to utilize technology. Since participants were to complete tasks on English-language website mockups, they had to be competent in English. The U.S.A.-based requirement was to fashion culturally-relevant data (and also to comply with the European Union and the People's Republic of China's stringent data privacy policies). And since the wireframes were all based on websites, limiting devices used to laptops or desktops was intended to make the experience feel more realistic, which was a primary goal of this study as discussed in Section \hyperlink{1.4}{1.4}.

    Prolific participants were compensated a base rate of \$1 USD for the roughly 5-minute-long task, from the adjusted \$12/hour Prolific recommended wage, as mentioned above, in Section \hyperlink{2.1}{2.1}, which would be provided after manually screening for abnormal results. Once participants were confirmed to meet study prerequisites, they were then redirected to Maze for the actual experiment.

\hypertarget{2.2.2}{\subsection{Maze and User Tests}}
    As mentioned above, Maze is a UX testing platform. This means that UI/UX professionals use this platform to test different design choices for their interfaces and record relevant data from participants. There are various platforms of this nature, including PlaybookUX, Userbrain, Lookback, etc., but Maze was chosen for this thesis pricing and recorded data types (including click heatmaps), amongst other features \cite{unmoderated_user-testing_tools-nng_nodate}. Maze was also used in a previous Reed Psychology thesis on UX design \cite{bec_psychology_2023}, which also informed this choice.

    In Maze, participants are once again screened for the attributes listed above in Section \hyperlink{2.2.1}{2.2.1}. They are additionally asked to input their Prolific IDs, which allows for compensation post-experiment. Then, after implying their consent, the meat of the study begins, starting with general study instructions. The instructions read as follows: \\

    \noindent\fbox{%
        \parbox{\textwidth}{You will be presented with 6 different websites. For each website, you will be given a task. Please click the correct feature to complete the task. In the case that tasks are complex and require 1+ steps, please click on the correct feature for the first step of the task. You are encouraged to complete these tasks as accurately and quickly as possible.\\

        Click continue to proceed.
        }%
    }\\

    Then, participants were randomly placed into one of three groups. Each group was exposed to a different variant of icon type ($Text-Only$, $Icon-Only$, $Icon-with-Text-Labels$); an example of each of these variants is shown below in Figure \hyperlink{3vars}{2.2}. All groups were exposed to the same six web UIs and tasks in the same order. Specifically, for each of the six tasks, the participants would first see a screen with the task on it, then would have to click to proceed to the web UI to complete the task. As explained in the instructional blurb above, completion of a task would be measured by when the participant clicked on the correct icon to complete the first step of a task. Creating UIs participants could interact with was achieved with a Figma prototype, which will be explained more in Section \hyperlink{2.2.3}{2.2.3} below. 
    
    \begin{figure}[H]
        \hypertarget{3vars}{}
        \begin{center}
            \subfigure[]{\includegraphics[width=.7\textwidth]{thesistemplate_2020-04-24/figures/mogeek IWTL.png}}
            \subfigure[]{\includegraphics[width=.7\textwidth]{thesistemplate_2020-04-24/figures/mogeek IO.png}}
            \subfigure[]{\includegraphics[width=.7\textwidth]{thesistemplate_2020-04-24/figures/mogeek TO.png}}

        \caption{Search engine wireframes on Figma with 3 variants: (a) Text-Only, (b) Icon-Only, and (c) Icon with Text Labels}
        \label{fig:placeholder}
        \end{center}
    \end{figure}
    
    In the study, the dependent variables (DVs) measured were Speed ((Time when correct feature is clicked) - (Time when started task)) and Misclick Rate ($\frac{\text{Number of incorrect clicks}}{\text{Number of total clicks}}$).

    \hypertarget{2.2.3}{\subsection{Figma and UI/UX Design Choices}}
        The six website types were chosen to mimic popular website functions, which was quantified through Cloudflare's Top 100 Domains list \cite{cloudflare}. The six chosen website types were: search engine (e.g. Google Search), e-commerce (e.g. Amazon), food ordering (e.g. GrubHub), hotel booking (e.g. Expedia), library (e.g. OverDrive), and video streaming (e.g. YouTube). In order to make realistic web UIs, the wireframes were modeled after mainly real Australian (thus, not frequented by Americans). The only two non-Aussie sites were Mojeek and Gom Player, which were selected nonetheless for their relative obscurity. Australian websites were chosen over non-English webpages to avoid translation errors. The chosen model webpages were \textit{Mojeek} for search engines \cite{mojeek}, \textit{Universal Store} for e-commerce \cite{universal_store}, \textit{Fishbowl} for food ordering \cite{fishbowl}, \textit{Vulcan Hotel} for hotel booking \cite{vulcan}, \textit{Penrith City Libraries} for library \cite{penrith_libraries}, and \textit{Gom Player} for video streaming \cite{gom_video_player}. The different model websites were also chosen to have a variety of perceived complexity, which was qualitatively assessed upon selection. Simply put, complexity was determined by the number of buttons on screen (which is akin to Grahame et al.\@'s definition of clutter \cite{grahame_age_2004}), the placement of said buttons, and the difficulty of the task paired with the website UI.

        All mockups were  created on Figma, an industry-standard design and prototyping tool used specifically for diagramming technological interfaces (or UIs). It has a wide range of capabilities, including prototyping, that allow UI/UX designers to not only mock-up their UIs, but also add interactive elements to these otherwise stagnant UIs \cite{what_is_figma}.
        
        All icons were from Google's \textit{Material Design Icon Library}, which contains over 2,500 modern, clean glyphs of all meanings \cite{material_library}. To use the terminology introduced in Sections \hyperlink{1.3.3}{1.3.3} and \hyperlink{1.3.4}{1.3.4}, the \textit{Material Design Icon Library} consists of flat icons of varying semantic distances (with some abstract and some concrete). The \textit{Material Design Icon Library} has integrations with Figma as well, allowing the icons to be altered in terms of dimension and color to match the aesthetics of these websites better. This library was also chosen for its previous use in Song et al. \cite{song_developing_2024}.  

        Most importantly, Figma has prototypes, which allow UI/UX designers to introduce interactive components to their wireframes. In order to create three separate participant groups per variant (text-only, icon-only, icon-with-text-labels), there were three prototyped user flows, or sequences (a.k.a. \enquote{flows}) in which users view the Figma mockups. The ordering of the websites was consistent between groups,the only thing that changed was the icon type on each site. As mentioned in Section \hyperlink{2.2.2}{2.2.2}, users from all groups first saw the instructions for tasks. Descriptions of tasks were centered on an otherwise blank Figma frame; when clicked, they would redirect the user to the corresponding web UI where the task would take place. Each participant group saw a different variant of the website. This was achieved by creating three of each task frame that would redirect to the correct variant upon user interaction, which you can see in Figures \hyperlink{figma_wideshot}{2.3} and \hyperlink{figma_zoom}{2.4}. All web UIs were coded such that the users were redirected to the next task frame when the correct feature corresponding to the task was clicked. Upon completing the sixth and last task, users were redirected to a blank frame. This blank frame was what Maze recognized as the end of the user flow, which would conclude the Figma prototype and allow users to continue through Maze. The Figma user flows were directly embedded into Maze via plug-in. 

        \begin{figure}[H]
            \hypertarget{figma_wideshot}{}
            \centering
            \includegraphics[width=1\linewidth]{Screenshot 2026-02-15 at 14.14.03.png}
            \caption{Entire Figma file with user flows}
            \label{fig:placeholder}
        \end{figure}

        \begin{figure}[H]
            \hypertarget{figma_zoom}{}
            \centering
            \includegraphics[width=1\linewidth]{zoom.png}
            \caption{\enquote{Search engine} section of Figma file with user flows}
            \label{fig:placeholder}
        \end{figure}

\hypertarget{2.3}{\section{Data Analysis}}
    \hypertarget{2.3.1}{\subsection{Data Recording and Migration}} 
        All data was recorded by Maze and migrated by hand into an Excel spreadsheet. The spreadsheet was formatted such that each participant had six entries, one per website. Each entry contained the following information: 
        \begin{enumerate}
            \item $Maze_ID$ (a string of integers)
            \item $Prolific_ID$ (a string of integers and characters)
            \item $Age_Group$ (either 18-35 or 60+)
            \item $Icon_Type$ (either Text Only, Icon Only, or Icon with Text Labels)
            \item $Website$ (an integer between 1-6 corresponding to the order of appearance of the website within the user flow)
            \item $Speed$ (floating point number, in seconds)
            \item $Misclick_Rate$ (floating point number, represented as a percentage of wrong clicks over total number of clicks)
        \end{enumerate}

        User responses were only migrated to Excel if they had fully completed the study, meaning that they found correct solutions to all six websites. Users who did not complete all six websites or dropped out before being assigned to an Icon Type condition (which is an important distinction for later in Section \hyperlink{2.3.3}{2.3.3}) were not included within the final Excel sheet.
        
        The spreadsheet ended up containing a total of 1440 rows, corresponding to 240 unique participants. In order to process the data in Python (see Section \hyperlink{2.3.2}{2.3.2}), the Excel spreadsheet was first converted into a \textit{.csv} file.

    \hypertarget{2.3.2}{\subsection{Tech Stack}}

        All data analysis was performed in Python, which enabled by various data-related Python libraries. The code for all of this study's data analysis, as well as the anonymized \textit{.csv} file mentioned in Section \hyperlink{2.3.1}{2.3.1}, can be found in the Github repository linked here: \url{https://github.com/00eemsy/thesis}.
    
        The primary Python library employed was \textit{pandas} (lowercase intentional), a popular, open-source data analysis library \cite{pandas}. This was used to access the data within the \textit{.csv} file, clean the data, and to perform data manipulation required for both the descriptive statistics and actual linear regression (see below in Sections \hyperlink{2.3.3}{2.3.3} and \hyperlink{2.3.4}{2.3.4}). For instance, \textit{pandas} was employed in the descriptive statistics to calculate the average speed of participants per website (see Figure \hyperlink{avg_speed}{3.1}).
    
        \textit{Matplotlib}, another popular, open-source Python library, was used instead for data visualization \cite{matplotlib}. All of the data visualizations found in this thesis, from bar charts, histograms, to scatterplots, were constructed by \textit{Matplotlib} with \textit{pandas} findings. 
    
        Besides these main two libraries, several other libraries were employed in one-off cases, including \textit{statsmodel} (lowercase intentional) for mixed-effects linear regression calculations \cite{statsmodel_mixedlm}, \textit{scipy} (lowercase intentional) for chi-squared calculations \cite{scipy_chisquare}, and \textit{numpy} (lowercase intentional) for array manipulation in tandem with \textit{pandas} calculations \cite{numpy}.

    \hypertarget{2.3.3}{\subsection{Descriptive Statistics}}
        Performing descriptive statistics analysis before the formal linear regression was necessary to achieve multiple goals. The main goal of these descriptive statistics, besides creating some casual familiarity with the data, was determining exclusion criteria. As mentioned above in Section \hyperlink{2.1}{2.1}, one of the biggest drawbacks to online survey platforms is noisy data. This can be caused by participants speeding through instructions \cite{palan_prolificacsubject_2018}, randomly clicking through UX tests, or even AI bot inaccuracies \cite{peterson_impact_nodate, asher_chatbots_2025, christoforou_generative_2024} (see more in Section \hyperlink{2.1.1}{2.1.1}). So cleaning up data is a priority to isolate the signal within the variance.

        The second goal of these descriptive statistics was making sense of the drop-off rate for those assigned an $Icon_Type$ variant. As mentioned in Section \hyperlink{2.3.1}{2.3.1}, this is related, but distinct, from those who dropped-off before being assigned a variant not assigned an Icon Type variant and total drop-off rate. Specifically, their relationship is as follows:

        $$\text{Total Drop-off} = \text{Drop-off}_{\text{Assigned Icon Type}} + \text{Drop-off}_{\text{Not Assigned Icon Type}} $$
        
        Maze recorded a total of 393 unique attempts (meaning there was an initial participant pool of size 393), but only 240 successfully completed attempts, was a cause for concern given the $\sim38.9\%$ total drop-off rate. To this end, descriptive statistics were supposed to shed light on what created this large total drop-off rate, and whether this had something to do either with the areas of interest (Age Group and Icon Type), which is why we look specifically at $\text{Drop-off}_{\text{Assigned Icon Type}}$.

        In terms of exclusion criteria, most of the work was looking at the dependent variables (DV; Speed, Misclick Rate) in relation to the various independent variables (IV; Age Group, Website, Icon Type, etc.). Average Speed was considered both per Website and summed across all six websites. To help contextualize these calculations, standard deviation was also calculated. More detailed numbers will be discussed in Section \hyperlink{3}{3} below. 

        This data was then plotted in various manners to help visualize the information. The three main vehicles to visualize this data were bar charts (for visualizing IV/DV correlations), histograms (for plotting distributions of DVs), and scatterplots (to compare between DVs Speed and Misclick Rate). 

        As for $\text{Drop-off}_{\text{Assigned Icon Type}}$ calculations, the main route for analysis was a chi-square test of independence. A chi-square test of independence decides if two categorical variables are associated \cite{discovery_chi-square}. The first chi-square test was performed between the number of participants per Age Group and Icon Type who finished the experiment versus those who did not finish but were assigned a variant. The second chi-square test was performed twice, once per Age Group, between the number of participants per Icon Type who finished the experiment versus those who did not finish but were assigned a variant. To clarify, those who dropped off before making it to the websites were not included in this measure. 

    \hypertarget{2.3.4}{\subsection{Mixed Effects Linear Regression}}

    Many hypothesis testing methods were considered before settling on a mixed effects linear regression. First, a t-test, a simpler model of hypothesis testing, was considered. These are ideal for testing between two conditions, but given that there is a 5\% false positive error rate per test that compounds \cite{one-way}, comparing between more than two variants (given the three Icon Type variants) was not viable. This could be mitigated by combining scores between variants (for instance, comparing $IWTL - IO$ versus $TO - IO$), but this was not as straightforward and therefore not pursued.

    ANOVAs were also considered, mainly because of their frequent appearance in the literature \cite{ganor_designing_nodate, schroder_effects_2008, liu_study_2023, dosso_how_2021, wu_how_2022, xi_influence_2018, song_developing_2024, grahame_age_2004, yu_reducing_2024}. And an ANOVA is mathematically identical to a linear regression as a special type of linear model. But this is confounded by the introduction of the Participant random effect. Since each participant brings baseline variations in speeds and accuracies, and each participant completes six tasks across the six websites, observations are not independent of one another due to random variation between participants. Classical ANOVAs assume independence between observations \cite{one_way_anova}, and repeated-measures ANOVAs, despite allowing for non-independence between observations, has less flexibility than necessary to process this data (for instance, it cannot handle unbalanced data well, which this table asserts is the case with for dataset) \cite{muhammad_guidelines_nodate}. On the other hand, mixed effects linear regressions are naturally able to handle this co-dependence with a random effect, which below is represented as $u_i$.

    Finally, the exact formulation of the linear regression has changed over time. Originally, the formula was:
    $$Outcome\sim Age + Icon + Website +(Age * Icon)$$ 
    
    in R notation and 
    
    \begin{gather*}
        DV = \beta_0 + \beta_1(Age_{60+})+(\beta_2(Icon_{IWTL})+\beta_3(Icon_{TO}))+\\(\beta_4(Website_2)+\beta_5(Website_3)+\beta_6(Website_4) + \beta_7(Website_5) + \beta_8(Website_6)) + \epsilon
    \end{gather*}

    when written out in full.
    
    This equation measures DVs (either Speed or Misclick Rate) from the main effects of Age Group, Icon Type, Website, and the interaction term between Age Group and Icon Type. However, this misses a key random effect, namely that of Participant, since baseline variance between Participants cannot be controlled for.

    The final equation required a mix of both fixed and random effect variables, hence why this thesis employs a mixed effects linear regression for the bulk of analysis. In R-speak, this equation is 
    
    $$Outcome\sim Age + Icon + Website +(Age * Icon) + (1|Participant)$$
    
    This is equivalent to

    \begin{gather*}
        DV_{ij} = \beta_0 + \beta_1(Age_{60+})_i+(\beta_2(Icon_{IWTL})+\beta_3(Icon_{TO}))_i+\\(\beta_4(Website_2)+\beta_5(Website_3)+\beta_6(Website_4) + \beta_7(Website_5) + \beta_8(Website_6))_j + u_i + \epsilon_{ij}
    \end{gather*}

    This is later amended to the following after applying exclusion criteria (which is explained more in Section \hyperlink{2.3.4}{2.3.4}):\\
    \begin{center}
        \fbox{\parbox{12.5cm}{$DV_{ij} = \beta_0 + \beta_1(Age_{60+})_i+(\beta_2(Icon_{IWTL})+\beta_3(Icon_{TO}))_i+\\(\beta_4(Website_4) + \beta_5(Website_5) + \beta_6(Website_6))_j + u_i + \epsilon_{ij}$}}\\
    \end{center}
    \bigbreak
    
    As mentioned previously, $u_i$ represents the Participant random effect. The $i$ indexes individual participants, whereas $j$ indexes websites, meaning $i\in[1,240]$ and $j\in [1,6]$. The indices correlate to the two levels of observation, namely the Website-level observations ($j$) (this includes all Website terms) and the Participants-level ($i$) (this includes Age, Icon Type, and $u_i$).

    With so many terms, overfitting was a risk. Overfitting occurs when terms are added to a model to the point that the model no longer represents the generalizable phenomenon and instead models the specific dataset instead \cite{perraillon}. In order to avoid this, an additive approach was employed, meaning that terms was added onto the model one-by-one and model iterations were compared at each step for relative fitness. Model 0 was the intercept-only model, model 1 added in Age Group, model 2 added in Icon Type, model 3 added in Website, and finally, model 4 added in the Age Group/Icon Type interaction term. Furthermore, Akaike Information Criterion, the performance metric for regression models employed also further negates this risk. AIC for linear regressions is:

    $$AIC_p = n\times \ln\left(\frac{SSE_p}{n}\right)+2p$$

    The parameters of AIC are as follows: $p$ is the number of estimated parameters, $n$ is the number of observations, and $SSE$ is the residual sum of squares. Essentially, this statistic estimates the amount of information lost when using a certain model \cite{aic}. One AIC's magnitude on its own does not rela any information about goodness of fit, it is when it's imployed comparatively to other models' AICs where the $\Delta AIC = AIC_{i-1}-AIC_{i}$ matters. That being said, the lowest AIC retained the most information, meaning it was the more accurate out of the compared models, meaning the term added provided substance to the model's predictive power without overfitting the data.

    There are several metrics for comparing performance between regression models, including $R^2$, but $AIC$ was the best choice because it penalizes overfitting and balances goodness of fit and parsimony well \cite{perraillon}. Specifically, the $+2p$ penalizes having a larger number of parameters \cite{perraillon}, which makes it better for combating overfitting as compared to other statistics such as $R^2$.
    
    When calculating AIC, several pairs of models were tested. The difference in AIC between models 0 and 1 was looked at to assess the impact of Age Group on the model. Models 1 and 2 were compared to estimate Icon Type's influence, models 2 and 3 estimated Website's influence, and finally, the $\Delta AIC$ of models 3 and 4 assess the importance of the Age Group and Icon Type interaction term. Most important for our study was seeing how the interaction between Age Group and Icon Type affected the model's fitness. 

    After deriving the most fit model, the significance levels of various variants were derived in comparison to \enquote{baseline} conditions (for instance, $IO$ was chosen as the baseline for Icon Type, meaning $IWTL$ and $TO$ were compared against $IO$'s results). 
 
    This Mixed Effects Linear Model approach was employed on a plethora of exclusion criteria, which were determined via descriptive statistics, which are discussed in Section \hyperlink{2.3.3}{2.3.3} and \hyperlink{3.1}{3.1}. Tests were done with the following exclusions:
    
    \begin{enumerate}
        \item None
        \item Websites 1 and 3
        \item Average Misclick Rate $>80\%$
        \item Total Speed $>200$s
        \item Total Speed $>200$s or Average Misclick Rate $>80\%$
    \end{enumerate}

    The second exclusion criteria, excluding all results from websites 1 and 3 entirely, was tested because websites 1 and 3 consistently performed worse in terms of both DVs (which will be discussed more thoroughly in Section \hyperlink{3}{3}). To conjecture, website 3 was a complex, visually busy UI (see Figure \hyperlink{w3}{2.5}) with a difficult task that relied on clicking a small button (\enquote{Remove Mixed Greens from Spicy Salmon order}), making for worse results across both DVs. And website 1, despite having a relatively simpler UI due to being a search engine (see Figure \hyperlink{3vars}{2.2}), had comparable levels of error possibly due to its appearance as first in the sequence of websites for all users. This user behavior was discovered during an informal, one-man pilot study completed by advisor Dr.\@ Adam Groce, where he clicked too fast on the first screen and skipped past vital instructions for completion of Website 1's task. If PhD doctors are not immune to this, it is probably a reasonable conjecture to generalize this phenomenon to the general populace of participants.
        
    \begin{figure}[H]
        \hypertarget{w3}{}
        \centering
        \includegraphics[width=1\linewidth]{w3.png}
        \caption{Website 3 UI (Text Only variant)}
    \end{figure}

    Figure \hyperlink{tsvamr}{2.6} was the initial justification for utilizing average Misclick Rate and total Speed as potential exclusion criteria. Data points usually fell into the lower left-hand corner, signifying faster and more accurate responses. However, a couple scatter far beyond that larger cluster, which was what created the $>200$s total Speed/$>80\%$ average Misclick Rate joint exclusion criteria. The red dotted line in the plot visualizes that cut-off.

    \begin{figure}[H]
        \hypertarget{tsvamr}{}
        \centering
        \includegraphics[width=1\linewidth]{total_speed_vs_avg_mr_v2.png}
        \caption{Total Speed vs average Misclick Rate by Age Group}
    \end{figure}

    \begin{figure}[H]
        \hyperlink{tsvmmr}{}
        \centering
        \includegraphics[width=1\linewidth]{total_speed_vs_median_mr.png}
        \caption{Total Speed vs median Misclick Rate by Age Group}
    \end{figure}

    The use of total Speed and average Misclick Rate was intentional as well. Given the variation in speed per Website, total Speed (the sum of all six speeds across the six different UIs per user) provided a more holistic interpretation of a user's performance. Average Misclick Rate (the average Misclick Rate across all six tasks per user) was also not immune to that variance, and since we were aggregating Speed by participant, there were only a couple of measures we could use to represent per-user Misclick Rate. 
    
    As a side note, Misclick Rate is calculated as:
    
    $$\frac{\text{Number of incorrect clicks}}{\text{Number of total clicks}}$$
    \bigbreak
    This fraction initially grows fast, but then plateaus early on, as you can see below in Figure \hyperlink{mr_growth}{2.8}, resulting in \enquote{diminishing returns}. This allows for easily skewed mean calculations if a participant found one website hard but breezed through the rest. 
    
    \begin{figure}[H]
        \hypertarget{mr_growth}{}
            \centering
            \begin{tabular}{| p{3.5cm} | p{3.5cm} |}
                \hline
                \# of Wrong Clicks & Misclick Rate (\%)\\
                \hline
                0 & 0\\
                \hline
                1 & 50.000\\
                \hline
                2 & 66.667\\
                \hline
                3 & 75.000\\
                \hline
                4 & 80.000\\
                \hline
                5 & 83.334\\
                \hline
                6 & 85.714\\
                \hline
                7 & 87.500\\
                \hline
                8 & 88.889\\
                \hline
                9 & 90.000\\
                \hline
                10 & 90.909\\
                \hline
            \end{tabular}
            \caption{Misclick Rate Growth}
        \label{fig:placeholder}
    \end{figure}
    
    Thus, median Misclick Rate (median Misclick Rate across all six UIs per user) was also considered. However, Figure \hyperlink{tsvmmr}{2.7} was less clustered within a singular location and therefore less easily segregated. Hence, average Misclick Rate was chosen over median Misclick Rate. So exclusion criteria 3-5 were all extracted from Figure \hyperlink{tsvamr}{2.6}. 

    Users with either a total Speed over 200s or an average Misclick Rate of over 80\% were completely excluded, even if some of their per-Website responses had Misclick Rates under $80\%$.

    In the end, the second exclusion criteria of no Websites 1 and 3 was chosen. Reasoning justifying this choice will be presented in Section \hyperlink{3.1.1}{3.1.1}, but the reasoning as to why not exclusion criteria 3-5 will be presented now.

    The main reason these exclusion criteria were not adopted was the lack of contextual justification for the thresholds, unlike the exclusion of Websites 1 and 3, which was based on identifiable interface characteristics. Excluding users based on average Speed or Misclick Rate required selecting arbitrary cut-off values that were not grounded in the task design or user behavior.

    In addition, these thresholds removed only eight participants and did not substantially reduce variability in the dataset. This suggests that the observed noise was not driven solely by a small number of extreme users. Instead, much of the variability appeared at the website level: many participants performed consistently on most websites but had unusually poor performance on one or two tasks. Because the exclusion criteria were based on participant-level averages, these noisy observations remained in the dataset even when the participant was retained. As a result, participant-level exclusions based on average Speed or Misclick Rate were not effective in meaningfully improving data quality.

    % The main reasoning for not using these exclusions is no proper contextual justification, unlike the exclusion of Websites 1 and 3. But alongside this, only excluding eight users did not cut down enough noise. Although Figure \hyperlink{tsvamr}{2.6} seems like a reasonable cut-off that excludes a lot of noise, the reality is that this excludes a total of 8 users from the data. But with a population of over 200 users and the amount of noise created, removing 8 users's results was simply not enough to clean the data to make it usable. For reference, the standard deviation for Speed is 16.683 and the standard deviation for Misclick Rate is 37.172-- standard deviation is variance, which includes both variation due to noise and variation due to signal, but these numbers are particularly high. For instance, many users had good performance on all but one or two Websites, making their average Misclick Rate and total Speed less than the cut-off thresholds but still noisy nonetheless. 
    
\hypertarget{3}{\chapter{Results}}

    This chapter will be split into two sections mirroring the split between Sections \hyperlink{2.3.3}{2.3.3}, which was about descriptive statistics, and \hyperlink{2.3.4}{2.3.4}, which was about mixed effects linear regressions. Each section will give a thorough summary of quantitative findings, and complete findings can be found within the Appendix. 

    \hypertarget{3.1}{\section{Descriptive Statistics}}
    
    Descriptive statistics, as mentioned previously in Section \hyperlink{2.3.3}{2.3.3}, were performed before the formal linear regression to provide context on a couple of fronts, namely exclusion criteria and drop-off rate. Descriptive statistics were run on the full set of data without exclusions. The subsequent subsections are organized in lieu of these initial interests, with additional important statistics contextualizing the findings of the linear regression sprinkled in between.

    \hypertarget{3.1.1}{\subsection{Exclusion Criteria}}

    The final exclusion criteria chosen from those listed in Section \hyperlink{2.3.4}{2.3.4} was the exclusions of Websites 1 and 3 entirely. That meant that the mixed-effects linear regression was performed solely from the data recorded from Websites 2, 4, 5, and 6. This exclusion criteria shaved the initial dataset of 1440 entries down to 960 entires, or exactly two-thirds of the original amount. 

    The goal of this subsection is to provide justification, both quantitative and qualitative, for this exclusion criteria. A high-level overview of the argument for this criteria can be found in Section \hyperlink{2.3.4}{2.3.4}.

    The first statistics helpful in visualizing the validity of this exclusion criteria are Tables \hyperlink{avg_speed}{3.1} and \hyperlink{avg_mr}{3.2}. 

    \begin{table}[H]
            \hypertarget{avg_speed}{}
            \centering
            \begin{tabular}{| p{3.5cm} | p{3.5cm} |}
                \hline
                Website & Speed (s)\\
                \hline
                1 & 17.155\\
                \hline
                2 & 7.858\\
                \hline
                3 & 23.814\\
                \hline
                4 & 7.0942\\
                \hline
                5 & 6.434\\
                \hline
                6 & 5.506\\
                \hline
            \end{tabular}
            \caption{Average Speed per Website}
            \label{fig:placeholder}
        \end{table}

    \begin{table}[H]
            \hypertarget{avg_mr}{}
            \centering
            \begin{tabular}{| p{3.5cm} | p{3.5cm} |}
                \hline
                Website & Misclick Rate (\%)\\
                \hline
                1 & 42.978\\
                \hline
                2 & 13.680\\
                \hline
                3 & 64.688\\
                \hline
                4 & 22.933\\
                \hline
                5 & 10.863\\
                \hline
                6 &  9.601\\
                \hline
            \end{tabular}
            \caption{Average Misclick Rate per Website}
            \label{fig:placeholder}
        \end{table}
    
    In both DVs, Websites 1 and 3 perform notably worse. Of note is the fact that Website 3 has both worse Speed and Misclick Rate than Website 1 (23.814s v.s. 17.155s, 64.688\% v.s. 42.978\%). 
    % This backs the original conjecture that Website 3's difficulty was rooted in its UI's complexity and that Website 1's was rooted in the order of appearance. Although some participants, like Dr. Adam Groce (see Section \hyperlink{2.3.4}{2.3.4}), perhaps still clicked too fast through Website 1's instructions, this was not a universal phenomenon. This discrepancy points to Website 1's UI not being an issue, and that those who did proceed with caution did not find major difficulty with the website otherwise. However, Website 3's complexity trended more towards universal than otherwise. This can also be seen across Age Groups in tables \hyperlink{avg_speed_age}{Average Speed per Website and Age Group} and \hyperlink{avg_mr_age}{Average Misclick Rate per Website and Age Group}.

    \begin{figure}[H]
        \hypertarget{avg_speed_age}{}
        \centering
        \begin{tabular}{| p{3.5cm} | p{3.5cm} | p{3.5cm}|}
            \hline
            Website & Age Group (yr) & Speed (s)\\
            \hline
            1 & 18-35 & 9.126983\\
            \hline
            & 60+ & 24.665565\\
            \hline
            2 & 18-35 & 4.263707\\
            \hline
            & 60+ & 11.219919\\
            \hline
            3 & 18-35 & 15.096724\\
            \hline
            & 60+ & 31.968065\\
            \hline
            4 & 18-35 &  4.294655\\
            \hline
            & 60+ & 9.713145\\
            \hline
            5 & 18-35 &  4.368362\\
            \hline
            & 60+ & 8.365887\\
            \hline
            6 & 18-35 &  3.781466\\
            \hline
            & 60+ & 7.118629\\
            \hline
        \end{tabular}
        \caption{Average Speed per Website and Age Group}
        \label{fig:placeholder}
    \end{figure}

    \begin{figure}[H]
        \hypertarget{avg_mr_age}{}
        \centering
        \begin{tabular}{| p{3.5cm} | p{3.5cm} | p{3.5cm}|}
            \hline
            Website & Age Group (yr) & Misclick Rate (\%)\\
            \hline
            1 & 18-35 & 37.704310\\
            \hline
            & 60+ & 47.910484\\
            \hline
            2 & 18-35 & 7.801724\\
            \hline
            & 60+ & 19.179032\\
            \hline
            3 & 18-35 & 61.812931\\
            \hline
            & 60+ & 67.377419\\
            \hline
            4 & 18-35 &  20.285345\\
            \hline
            & 60+ & 25.408871\\
            \hline
            5 & 18-35 &  4.120690\\
            \hline
            & 60+ & 17.169355\\
            \hline
            6 & 18-35 &  6.726724\\
            \hline
            & 60+ & 12.289516\\
            \hline
        \end{tabular}
        \caption{Average Misclick Rate per Website and Age Group}
        \label{fig:placeholder}
    \end{figure}
    
    No one Age Group finds disproportionate amounts of trouble with Websites 1 and 3 as compared to one another, they are both negatively affected. This points to experimental design error (and noise) as opposed to any significant age-related effect (signal). 

    Variance is also important in contextualizing these findings. Given the hardships posed by Websites 1 and 3, the results reaped from each website happen to be noisier and therefore less useful. This can be seen in Figures \hyperlink{speed_var}{3.3} and \hyperlink{mr_var}{3.4}. 

    \begin{figure}[H]
        \hypertarget{speed_var}{}
        \centering
        \begin{tabular}{| p{3.5cm} | p{3.5cm} |}
            \hline
            Website & Standard Deviation\\
            \hline
            1 & 27.160808\\
            \hline
            2 & 8.896048\\
            \hline
            3 & 20.878590\\
            \hline
            4 & 7.135238\\
            \hline
            5 & 8.078409\\
            \hline
            6 &  5.290130\\
            \hline
        \end{tabular}
        \caption{Speed Standard Deviation per Website}
        \label{fig:placeholder}
    \end{figure}

    \begin{figure}[H]
        \hypertarget{mr_var}{}
        \centering
        \begin{tabular}{| p{3.5cm} | p{3.5cm} |}
            \hline
            Website & Standard Deviation\\
            \hline
            1 & 38.857985\\
            \hline
            2 & 28.404598\\
            \hline
            3 & 34.398600\\
            \hline
            4 & 33.121365\\
            \hline
            5 & 25.779740\\
            \hline
            6 & 24.817584\\
            \hline
        \end{tabular}
        \caption{Misclick Rate Standard Deviation per Website}
        \label{fig:placeholder}
    \end{figure}

    In Figure \hyperlink{speed_var}{3.3}, it is clearer to see the greater variances in Websites 1 (27.160808, 38.857985) and 3 (20.878590, 34.398600) versus the rest. Interesting to note is that the standard deviation on Website 4 is nearly that of Website 3's. 
    
    Some visualizations may help contextualize this variance. First, consider Figures \hyperlink{mr_distr_w1}{3.5} and \hyperlink{mr_distr_w3}{3.6}. After just two erroneous clicks, Misclick Rate shoots over 50\% (see Figure \hyperlink{mr_growth}{2.8}), creating two poles. Notice that in \hyperlink{mr_distr_w1}{Misclick Rate Distribution for Website 1} the ratio between the number of participants at either end are roughly the same, and that in \hyperlink{mr_distr_w3}{Misclick Rate Distribution for Website 3} there are more high Misclick Rates than low. 

    \begin{figure}[H]
        \hypertarget{mr_distr_w1}{}
        \centering
        \includegraphics[width=0.75\linewidth]{mr_distr_w1.png}
        \caption{Misclick Rate Distribution for Website 1}
    \end{figure}

    \begin{figure}[H]
        \hypertarget{mr_distr_w3}{}
        \centering
        \includegraphics[width=0.75\linewidth]{mr_distr_w3.png}
        \caption{Misclick Rate Distribution for Website 3}
    \end{figure}

    The distribution for Website 6 (shown in \hyperlink{mr_distr_w6}{Misclick Rate Distribution for Website 6}) also creates two tails, which is expected given the nature of Misclick Rates (anything between 0-50\% is mathematically impossible), but the amount of high Misclick Rates pales in comparison to Websites 1 and 3. 

    \begin{figure}[H]
        \hypertarget{mr_distr_w6}{}
        \centering
        \includegraphics[width=0.75\linewidth]{mr_distr_w6.png}
        \caption{Misclick Rate Distribution for Website 6}
    \end{figure}
      
    In sum, this points to noisier data from Websites 1 and 3, another case for excluding them from data analysis.

    \hypertarget{3.1.2}{\subsection{Text Only and 60+}}

    In this next section, descriptive statistics pointing towards the findings within my linear regression are provided (see Section \hyperlink{3.2}{3.2}). Namely, this data will show that there is a significant relationship between the Text Only (TO) Icon Type and the 60+ Age Group.

    As mentioned above in the Background Literature (see Section \hyperlink{1.3}{1.3}), there is debate as to whether Icon Only (IO) versus Text Only (TO) variants produce greater performance effects. The following two tables, \hyperlink{speed_per_icon}{Average Speed per Icon Type} and \hyperlink{mr_per_icon}{Average Misclick Rate per Icon Type}, exhibit a neck-and-neck similarity between IO and TO as well.

    \begin{figure}[H]
        \hypertarget{speed_per_icon}{}
        \centering
        \begin{tabular}{| p{5cm} | p{3.5cm} |}
            \hline
            Icon Type & Speed (s)\\
            \hline
            Icon Only & 10.869219\\
            \hline
            Icon with Text Labels & 12.367133\\
            \hline
            Text Only & 10.793081\\
            \hline
        \end{tabular}
        \caption{Average Speed per Icon Type}
        \label{fig:placeholder}
    \end{figure}

    \begin{figure}[H]
        \hypertarget{mr_per_icon}{}
        \centering
        \begin{tabular}{| p{5cm} | p{3.5cm} |}
            \hline
            Icon Type & Misclick Rate (\%)\\
            \hline
            Icon Only & 27.064135\\
            \hline
            Icon with Text Labels & 31.788667\\
            \hline
            Text Only & 24.039922\\
            \hline
        \end{tabular}
        \caption{Average Misclick Rate per Icon Type}
        \label{fig:placeholder}
    \end{figure}

    Icon with Text Labels (IWTL) performs worse across both DVs (12.367133s, 31.788667\%) than TO and IO, but interestingly enough, TO ekes ahead (10.793081s, 24.039922\%) as compared to IO (10.869219s, 27.064135\%) in both as well. Another interesting phenomenon arises when the data is further broken down by Age Group (see Figures \hyperlink{speed_per_icon_age}{3.10} and \hyperlink{mr_per_icon_age}{3.11}).

    \begin{figure}[H]
        \hypertarget{speed_per_icon_age}{}
        \centering
        \begin{tabular}{| p{5cm} | p{3.5cm} | p{3.5cm}|}
            \hline
            Icon Type & Age Group (yr) & Speed (s)\\
            \hline
            Icon Only & 18-35 & 5.954573\\
            \hline
            & 60+ & 15.661000\\
            \hline
            Icon with Text Labels & 18-35 & 8.061842\\
            \hline
            & 60+ & 16.788784\\
            \hline
            Text Only & 18-35 & 6.481325\\
            \hline
             & 60+ & 14.370922\\
            \hline
        \end{tabular}
        \caption{Average Speed per Icon Type and Age Group}
        \label{fig:placeholder}
    \end{figure}

    \begin{figure}[H]
        \hypertarget{mr_per_icon_age}{}
        \centering
        \begin{tabular}{| p{5cm} | p{3.5cm} | p{3.5cm}|}
            \hline
            Icon Type & Age Group (yr) & Speed (s)\\
            \hline
            Icon Only & 18-35 & 23.446581\\
            \hline
            & 60+ & 30.591250\\
            \hline
            Icon with Text Labels & 18-35 & 26.070175\\
            \hline
            & 60+ & 37.661712\\
            \hline
            Text Only & 18-35 & 19.785897\\
            \hline
             & 60+ & 27.569858\\
            \hline
        \end{tabular}
        \caption{Average Misclick Rate per Icon Type and Age Group}
        \label{fig:placeholder}
    \end{figure}

    Excitingly, the \enquote{best} Icon Type for average Speed between ages differs. The 18-35 Age Group has a faster performance on IO (5.954573s) than TO (6.481325s), but the 60+ Age Group has the opposite, with an average speed of 14.370922s on the TO condition and 15.661000s on the IO one. Misclick Rate does not exhibit this difference, with both Age Groups erroring less on TO versus IO, but this discrepancy between Age Groups points towards something larger, which is validated in the mixed effects linear regression later in Section \hyperlink{3.2}{3.2}. 

    \hypertarget{3.1.3}{\subsection{Drop-Off Rate}}

    As mentioned above in Section \hyperlink{2.3.3}{2.3.3}, comprehending the large total drop-off rate ($\sim38.9\%$) was the second goal of descriptive statistics. Specifically, the question was whether any of our variables (namely, Age Group and Icon Type) had any impact on drop-off rate, specifically the drop-off rate of those assigned Icon Type variants, which is addressed by Figure \hyperlink{finish}{3.12}.

    \begin{figure}[H]
        \hypertarget{finish}{}
        \centering
        \begin{tabular}{| p{4cm} | p{3cm} | p{2.5cm}| p{3cm} |}
            \hline
            Icon Type & Age Group (yr) & Finishers & Non-Finishers\\
            \hline
            Icon Only & 18-35 & 39 & 6\\
            \hline
            & 60+ & 40 & 17\\
            \hline
            Icon with Text Labels & 18-35 & 38 & 6\\
            \hline
            & 60+ & 37 & 23\\
            \hline
            Text Only & 18-35 & 39 & 4 \\
            \hline
             & 60+ & 47 & 20\\
            \hline
        \end{tabular}
        \caption{Finishers versus Non-Finishers by Icon Type and Age Group}
        \label{fig:placeholder}
    \end{figure}

    The number of finishing TO 60+ users (47) is high. This could be in relation to the number of drop-offs with Icon Type assignments versus drop-offs without Icon Type assignments, but a lot of this is obscured away by Maze, the data collection platform. 
    
    As mentioned in section \hyperlink{2.3.3}{2.3.3}, a chi-squared test of independence was run to see if the ratios of finishers versus non-finishers were similar across Icon Types and Age Groups. The returned p-value is $\approx0.0015$, which makes it significantly different. Most of it is most likely rooted in Age Group disparity. This makes sense given that decreased ability, both physically and intellectually, is caused in part by increased age (see Section \hyperlink{1.2}{1.2}). 

    To be completely sure, another round of chi-squared tests of independence were run between finishers to non-finishers, one with only the 18-35 Age Group and another with only the 60+ Age Group. Both tests resulted in nonsignificant p-values (0.7875462343821208 for 18-35, 0.5148579660895773 for 60+), which points to the fact that the significant difference was due to differences between Age Group and not Icon Type. But it could also be the case that the nonsignificant p-values could have been due to there not being enough data per condition to draw firm conclusions.

\hypertarget{3.2}{\section{Mixed Effects Linear Regression}}
    The exclusion criteria settled upon was excluding results from Websites 1 and 3 entirely, as justified in Section \hyperlink{3.1.1}{3.1.1}. The following linear regression results were calculated using the dataset with exclusion criteria applied. 

    As aforementioned in Section \hyperlink{2.3.4}{2.3.4}, linear regression models were computed in an additive manner, meaning that variables were inserted into the equation one by one to see whether they made a marginal impact on the fitness of the Linear Regression. Figure \hyperlink{models}{3.13} formally defines what each model version measures in R-speak. Note that each model is grouped by Participant, hence the Participant random effect (written as $(1 | Maze\text{ }ID)$).
    
    \begin{figure}[H]
        \hypertarget{models}{}
        \centering
        \begin{tabular}{| p{1.5cm} | p{13cm} |}
            \hline
            Model & Definition\\
            \hline
            Model 0 & $Speed\sim1 + (1 | Maze\text{ }ID)$\\
            \hline
            Model 1 & $Speed\sim Age\text{ }Group + (1 | Maze\text{ }ID)$\\
            \hline
            Model 2 & $Speed\sim Age\text{ }Group + Icon\text{ }Type+(1 | Maze\text{ }ID)$\\
            \hline
            Model 3 & $Speed\sim Age\text{ }Group + Icon\text{ }Type + Website+(1 | Maze\text{ }ID)$\\
            \hline
            Model 4 & $Speed\sim Age\text{ }Group + Icon\text{ }Type + Website + (Age\text{ }Group * Icon\text{ }Type)+(1 | Maze\text{ }ID)$\\
            \hline
        \end{tabular}
        \caption{Definitions of linear regression models}
        \label{fig:placeholder}
    \end{figure}
    
    This impact was measured with an Akaike Information Criterion (AIC) to determine the model of best fit. The model with the lowest AIC comparatively fits the best.

    \hypertarget{3.2.1}{\subsection{Speed}}

    The AIC for Speed is shown below in Figure \hyperlink{aic_speed}{3.14}:

    \begin{figure}[H]
        \hypertarget{aic_speed}{}
        \centering
        \begin{tabular}{| p{3.5cm} | p{7.5cm} |}
            \hline
            Model & Akaike Information Criterion\\
            \hline
            Model 0 & 6522.543660965796\\
            \hline
            Model 1 & 6457.99778589942\\
            \hline
            Model 2 & 6438.843481013769\\
            \hline
            Model 3 & 6427.435013923644\\
            \hline
            Model 4 & 6422.526364301187\\
            \hline
        \end{tabular}
        \caption{AIC for Models on Speed}
        \label{fig:placeholder}
    \end{figure}

    Model 4, or the full model of $Speed\sim Age\text{ }Group + Icon\text{ }Type + Website + (Age\text{ }Group \times Icon\text{ }Type)+(1|Maze\text{ }ID)$, performed the best. 

    To break it down even further, the $\Delta AIC$ between two models with only one term different between them implies the difference in support of this model over another, and therefore how much of the variance in data is explained by adding this term to the preexisting model \cite{burnham_information_2002}. Namely, for $\Delta AIC=(\text{Model with less terms}) - (\text{Model with more terms})$, $0<\Delta AIC < 2$ means the models have equivalent levels of support, $4<\Delta AIC <7$ means the model with the added term has some support for being a better fitted model, and $\Delta AIC >10$ means that the model with the added term has substantially better support for being a model of better fit. This not mean that in general these variables hold significance in respect to the DV, it just means that the variable identifies some signal amidst variance between the two models being compared.

    So that being said, the variance explained by Age Group ($\Delta AIC=64.5458750664$), Icon Type ($\Delta AIC=19.1543048857$), and Website ($\Delta AIC=11.4084670901$) is substantial, making these terms important to include in the final model. The interaction term between Age Group and Icon Type also has some support for inclusion with a $\Delta AIC=4.90864962246$.

    With this in mind, we run model 4 (see Figure \hyperlink{no_w1_w3_m4}{3.15} and linear model below). Note that like linear regressions, not all variants are represented in Figure \hyperlink{no_w1_w3_m4}{3.15}. The coefficients provided are comparing that variant to a randomly selected baseline variant. Age Group's baseline is 18-35, Icon Type's is IO, Website's is Website 2, and Age Group and Icon Type's is 18-35 and IO. 

    \begin{figure}[H]
        \hypertarget{no_w1_w3_m4}{}
        \centering
        \begin{tabular}{| p{5cm} | p{3.5cm} | p{3.5cm} | p{3.5cm} |}
            \hline
            & Coefficient & z & P$>|\text{z}|$\\
            \hline
            Intercept & 5.050 & 6.790 & 0.000\\
            \hline
            Age[60+] & 5.951 & 6.491 & 0.000\\
            \hline
            Icon[IWTL] & 1.072 & 1.155 & 0.248\\
            \hline
            Icon[TO] & -0.264 & -0.287 & 0.774\\
            \hline
            Website[4] & -0.764 & -1.130 & 0.190\\
            \hline
            Website[5] & -1.424 & -2.444 & 0.015\\
            \hline
            Website[6] & -2.352 & -4.037 & 0.000\\
            \hline
            Age[60+]:Icon[IWTL] & 0.544 & 0.414 & 0.679\\
            \hline
            Age[60+]:Icon[TO] & -3.018 & -2.371 & 0.018\\
            \hline
            Group Variance & 6.417 &&\\
            Scale & 40.7389 &&\\
            \hline
        \end{tabular}
        \caption{Results of Speed Model 4}
        \label{fig:placeholder}
    \end{figure}

    The results shown in Figure \hyperlink{no_w1_w3_m4}{3.15} can be modeled as follows:

    \begin{gather*}
        Speed = 5.050 + (5.951(Age_{60+})) + (1.072(Icon_{IWTL}) - 0.264(Icon_{TO})) + \\(-0.764(Website_4) - 1.424(Website_5) - 2.352(Website_6)) + \\(0.544(Age_{60+}\times Icon_{IWTL}) - 3.018(Age_{60+}\times Icon_{TO})) + u_i + \epsilon_{ij}
    \end{gather*}

    with $Var(u_i)=6.147$ and $Var(\epsilon_{ij})=40.7389$.

    There are a couple of things to note from this. First, and most importantly, there is a significant interaction between the 60+ Age Group and the TO Icon Type (p-value = 0.018). This is not only significant, but also a positive effect that shaves off just over 3 seconds on Speed (coefficient=-3.018). Thus, the effect of TO on the 60+ Age Group, as compared to 18-35, is stronger in the beneficial direction. 

    Some other variables were significant as well. Age Group's significance was expected (p-value $\approx 0.000$), as noted previously in the literature (see Chapter \hyperlink{1}{1}), adding nearly 6 seconds to Speed (coefficient=5.951). This also puts our main finding into context-- for a user in the 60+ Age Group, viewing the TO variant nearly halves the negative RT effects of age.
    
    Interestingly enough, Websites 5 (p-value=0.015) and 6 (p-value$\approx0.000$) are also significantly different than Website 2. This might be explained with the same logic that resulted in the exclusion of Website 1. Namely, that learning effect kicks in and allows users to complete later Websites more efficiently due to repeated exposure. This is further backed by the fact that the Speed boost increases between Websites 4 (coefficient=-0.764) to 6 (coefficient=-2.352), and that only Websites 5 and 6 are significant as compared to Website 4 (p-value=0.190).

    Looking at Icon Types, some interesting trends arise. The different signs of the coefficients mirror what was shown in Section \hyperlink{3.1.2}{3.1.2}-- that is, that IWTL (coefficient=5.951) looks to perform worse than both IO and TO, and that TO somewhat performs better than IO (coefficient=-0.264). Of course, neither of these values are significant (p-values=0.248, 0.774), but this changes when looking at Misclick Rate instead of Speed. 

\hypertarget{3.2.2}{\subsection{Misclick Rate}}

    Again, we begin by calculating AIC (see Figure \hyperlink{aic_mr}{3.16}).

    \begin{figure}[H]
        \hypertarget{aic_mr}{}
        \centering
        \begin{tabular}{| p{3.5cm} | p{7.5cm} |}
            \hline
            Model & Akaike Information Criterion\\
            \hline
            Model 0 & 9121.182184951773\\
            \hline
            Model 1 & 9108.807437991436\\
            \hline
            Model 2 & 9083.299844112014\\
            \hline
            Model 3 & 9048.092978443245\\
            \hline
            Model 4 & 9050.568425915011\\
            \hline
        \end{tabular}
        \caption{AIC for Models on Misclick Rate}
        \label{fig:placeholder}
    \end{figure}

    Age Group, Icon Type, and Website all show substantial $\Delta AIC$s, with $\Delta AIC$s of 12.3747469603, 25.5075938794, and 35.2068656688 respectively. However, the interaction term between Age Group and Icon Type actually has negative $\Delta AIC$ ($-2.47544747177$), meaning including the additional term does not explain any of the variance and actually negatively impacts the model's fitness. However, it is on a much smaller magnitude than the other AIC changes between variables.

    This means that Model 3 is our model of best fit, which denies that the Age Group and Icon Type interaction term helps with model fitness. This points to that interaction not being significant, which is further validated in Figures(\hyperlink{no_w1_w3_m3_mr}{3.17}, \hyperlink{no_w1_w3_m4_mr}{3.18}):

    \begin{figure}[H]
        \hypertarget{no_w1_w3_m3_mr}{}
        \centering
        \begin{tabular}{| p{5cm} | p{3.5cm} | p{3.5cm} | p{3.5cm} |}
            \hline
            & Coefficient & z & P$>|\text{z}|$\\
            \hline
            Intercept & 9.903 & 3.856 & 0.000\\
            \hline
            Age[60+] & 9.318 & 4.340 & 0.000\\
            \hline
            Icon[IWTL] & 6.255 & 2.337 & 0.019\\
            \hline
            Icon[TO] & -8.351 & -3.225 & 0.001\\
            \hline
            Website[4] & 9.252 & 4.085 & 0.000\\
            \hline
            Website[5] & -2.817 & -1.244 & 0.214\\
            \hline
            Website[6] & -2.352 & -4.037 & 0.000\\
            \hline
            Group Variance & 121.745 &&\\
            \hline
            Scale & 615.718 &&\\
            \hline
        \end{tabular}
        \caption{Results of Misclick Rate Model 3}
        \label{fig:placeholder}
    \end{figure}

    \begin{figure}[H]
        \hypertarget{no_w1_w3_m4_mr}{}
        \centering
        \begin{tabular}{| p{5cm} | p{3.5cm} | p{3.5cm} | p{3.5cm} |}
            \hline
            & Coefficient & z & P$>|\text{z}|$\\
            \hline
            Intercept & 10.567 & 3.533 & 0.000\\
            \hline
            Age[60+] & 8.007 &  2.150 & 0.032\\
            \hline
            Icon[IWTL] & 3.678 & 0.975 & 0.330\\
            \hline
            Icon[TO] & -7.813 & -2.085 & 0.037\\
            \hline
            Website[4] & 9.253 & 4.085 & 0.000\\
            \hline
            Website[5] & -2.817 & -1.244 & 0.214\\
            \hline
            Website[6] & -4.079 & -1.801 & 0.000\\
            \hline
            Age[60+]:Icon[IWTL] & 5.190 & 0.972 & 0.331\\
            \hline
            Age[60+]:Icon[TO] & -0.886 & -0.171 & 0.864\\
            \hline
            Group Variance & 120.000 & &\\
            \hline
            Scale & 615.718&&\\
            \hline
        \end{tabular}
        \caption{Results of Misclick Rate Model 4}
        \label{fig:placeholder}
    \end{figure}

    Model 3, the best linear model out of the group, can also be represented as:

    \begin{gather*}
        Misclick\text{ }Rate = 9.903 + (9.318(Age_{60+})) + (6.255(Icon_{IWTL}) - 8.351(Icon_{TO})) + \\(9.252(Website_4) -2.817(Website_5) -4.079(Website_6)) + u_i + \epsilon_{ij}
    \end{gather*}
    \bigbreak
    with $Var(u_i)=121.745$ and $Var(\epsilon_{ij})=615.718$.
    
    As shown in model 4, neither interaction term was significant (p-values=0.331, 0.864). So unlike Speed, Age Groups do not have a significantly different Misclick Rate when provided different Icon Type variants. 

    But in model 3, all Age Groups and Icon Type coefficients are significant (p-values=0.000, 0.000, 0.019, 0.001). Specifically, 60+ has a marked increase in Misclick Rate versus 18-35 (coefficient=9.318), IWTL increases Misclick Rate as opposed to IO (coefficient=6.255), and TO decreases Misclick Rate as opposed to IO (coefficient=-8.351). 

    This means that separately, these variables have a significant impact on Misclick Rate. But when given certain Icon Types, the changes in Misclick Rate does not differ strongly between Age Groups the same way Speed does. Again, this could be due to Misclick Rate's \enquote{diminishing returns}. You also see the impacts of this growth rate in the magnitude of the coefficients, as compared to those in the Speed model.

    As for Websites, the only significant difference was between Websites 4 and 2 (p-value$\approx0.000$). Not only that, but Website 4's Misclick Rate actually worsens as compared to 2's (coefficients=9.252). This differs from Speed in a couple of senses-- first, Websites 5 and 6 were significant in Speed and Website 4 was not, whereas Website 4 is significant and Websites 5 and 5 are not for Misclick Rate in model 4. But second, Website 4 exhibited a negative RT coefficient whereas Website 4 is performing worse in terms of errors. Even harkening back to Figure \hyperlink{avg_mr}{3.2}, there is a small uptick in Misclick Rate on Website 4 (22.932500\%), albeit nowhere as large as Websites 1 (42.977500\%) and 3 (64.687917\%). To conjecture, this could be the result of Website 4 appearing right after Website 3, arguably the most difficult, complex, and therefore frustrating UI of the batch. If Website 3 caused participants to give up and begin clicking around randomly, this frustration could have carried over to Website 4 without them realizing how much more straightforward the task was. After all, Mahmud et al.\@ talks at length about how their older users were much more easily frustrated than their younger age groups, making them repeat unhelpful actions without course correction \cite{mahmud_learning_2020}.

\hypertarget{3.3}{\section{Robustness Checks}}

    In the following section, the results of the linear regression discussed above in Section \hyperlink{3.2}{3.2} are compared to regressions performed on similar, alternative exclusion criteria. 

    After settling upon the exclusion criteria of no Websites 1 and 3, the alternative criteria to confirm robustness were:
    \begin{enumerate}
        \item None
        \item (No websites 1 and 3) or ($>200$s total Speed or $>80$\% Misclick Rate)
        \item (No websites 1 and 3) or ($>200$s total Speed)
        \item (No websites 1 and 3) or ($>80$\% Misclick Rate)
    \end{enumerate}

    Criteria $2-4$ were all chosen because of the alternative exclusion criteria listed in Section \hyperlink{2.3.4}{2.3.4}. As a reminder, in the case of criterion 2, if either the user's total Speed was above 200 seconds or their average Misclick Rate was over 80\%, all of their results would be excluded in line with Figure \hyperlink{tsvamr}{2.6}. When split into criteria 3 and 4, this logic still holds-- if an user doesn't meet the threshold, all of their results are removed from analysis.
    
    Given the Figure \hyperlink{tsvamr}{2.6}, the original exclusion criteria 2-5 had some justification, although not as strong of justification as the chosen criteria. Therefore, there is still merit in studying the union of both exclusion criteria, since it's worth looking at whether the findings from Section \hyperlink{3.2}{3.2} hold under either defensible exclusion strategy.

    The other options for exclusion criteria mentioned in Section \hyperlink{2.3.4}{2.3.4} were also tried independently, but deemed not as relevant to this thesis. Other robustness checks were also performed, but also deemed not as relevant. Both of these findings, however, can still be found in the Appendix.

\hypertarget{3.3.1}{\subsection{No Exclusions}}
    Again, following the format of Section \hyperlink{3.2}{3.2}, AIC results will first be shown, followed by relevant linear regression results.

    First, we discuss Speed. Model 3's AIC score was the lowest at 11813.533, which already diverges from Section \hyperlink{3.2}{3.2}'s findings. The following are snippets from the results from models 3 and 4 (see Figures \hyperlink{ne_m3_speed}{3.19} and \hyperlink{ne_m4_speed}{3.20}) that are most relevant for comparison.

    \begin{figure}[H]
        \hypertarget{ne_m3_speed}{}
        \centering
        \begin{tabular}{| p{5cm} | p{3.5cm} | p{3.5cm} | p{3.5cm} |}
            \hline
            & Coefficient & z & P$>|\text{z}|$\\
            \hline
            Icon[IWTL] & 1.612 & 1.343 & 0.179\\
            \hline
            Icon[TO] & -0.428 & -0.369 & 0.712\\
            \hline
        \end{tabular}
        \caption{Results of Speed Model 3 with No Exclusions}
        \label{fig:placeholder}
    \end{figure}

    \begin{figure}[H]
        \hypertarget{ne_m4_speed}{}
        \centering
        \begin{tabular}{| p{5cm} | p{3.5cm} | p{3.5cm} | p{3.5cm} |}
            \hline
            & Coefficient & z & P$>|\text{z}|$\\
            \hline
            Icon[IWTL] & 2.107 & 1.234 & 0.217\\
            \hline
            Icon[TO] & 0.527 & 0.310 & 0.756\\
            \hline
            Age[60+]:Icon[IWTL] & -0.979 & -0.405 & 0.685\\
            \hline
            Age[60+]:Icon[TO] & -1.817 & -0.776 & 0.438\\
            \hline
        \end{tabular}
        \caption{Results of Speed Model 4 with No Exclusions}
        \label{fig:placeholder}
    \end{figure}

    In model 4, neither Icon Types nor the interaction term exhibited any significance. The coefficients for these terms are outliers from what we see throughout the main regression (see \hyperlink{3.2}{3.2}) and the rest of the robustness check results. Both IWTL and TO are positive (coefficients=2.107, 0527), meaning that overall both performed worse than IO. This seems to contradict calculations for Figure \hyperlink{avg_speed}{3.1}, but this could be due to model error. However, both interaction terms had negative coefficients (coefficients=-0.979, -1.817), implying better performance from the 60+ group than the 18-35 group in response to these specific Icon Types. This also seems to contradict findings presented in Chapter \hyperlink{3}{3}, but again this could be due to model error.
    
    Again, IWTL got the closest out of the four terms to significance (p-value=0.217). As another note, all other terms not listed (Age Group, Websites) exhibited p-values of 0.000.

    As for Misclick Rate, model 3 once again outperforms the rest with an AIC of 13927.396, meaning the interaction terms do not benefit model fitness marginally.

    \begin{figure}[H]
        \hypertarget{ne_m3_mr}{}
        \centering
        \begin{tabular}{| p{5cm} | p{3.5cm} | p{3.5cm} | p{3.5cm} |}
            \hline
            & Coefficient & z & P$>|\text{z}|$\\
            \hline
            Icon[IWTL] & 4.838 &  1.852 & 0.064\\
            \hline
            Icon[TO] & -3.376 & -1.337 & 0.181\\
            \hline
        \end{tabular}
        \caption{Results of Misclick Rate model 3 with no exclusions}
        \label{fig:placeholder}
    \end{figure}

    \begin{figure}[H]
        \hypertarget{ne_m4_mr}{}
        \centering
        \begin{tabular}{| p{5cm} | p{3.5cm} | p{3.5cm} | p{3.5cm} |}
            \hline
            & Coefficient & z & P$>|\text{z}|$\\
            \hline
            Age[60+] & 7.145 & 1.963 & 0.050\\
            \hline
            Icon[IWTL] & 2.624 & 0.712 & 0.477\\
            \hline
            Icon[TO] & -3.661 & -1.000 & 0.317\\
            \hline
            Age[60+]:Icon[IWTL] & 4.447 &  0.853 & 0.394\\
            \hline
            Age[60+]:Icon[TO] & 0.639 & 0.127 & 0.899\\
            \hline
        \end{tabular}
        \caption{Results of Misclick Rate model 4 with no exclusions}
        \label{fig:placeholder}
    \end{figure}

    In model 3, IWTL ekes towards significance, but still doesn't hit the mark (p-value=0.064). Again, all other coefficients had p-values of 0.000.  
    
    In model 4, coefficients behave more in line with expectations (see Section \hyperlink{3.2}{3.2}) than Speed (coefficients=2.624, -3.661, 4.447, 0.639), but still nothing of significance. Interesting to note is that Age is barely no longer significant (p-value=0.050), meaning this model did not notice significantly different error rates between the Age Groups, which defies previous findings (see Sections \hyperlink{1}{1} and \hyperlink{3.2}{3.2}). All Website terms had p-values of 0.000.

\hypertarget{3.3.2}{\subsection{(No websites 1 and 3) or ($>200$s total Speed and or $>80$\% average Misclick Rate)}}

    Once again, we look at Speed first. Model 4 had the lowest comparative AIC (5271.914), which is in line with the findings of the chosen linear regression.

    \begin{figure}[H]
        \hypertarget{n13_8_ppl_speed_m4}{}
        \centering
        \begin{tabular}{| p{5cm} | p{3.5cm} | p{3.5cm} | p{3.5cm} |}
            \hline
            & Coefficient & z & P$>|\text{z}|$\\
            \hline
            Icon[IWTL] & 0.692 &  1.188 & 0.235\\
            \hline
            Icon[TO] & 0.026 & 0.045 & 0.964\\
            \hline
            Age[60+]:Icon[IWTL] & 0.253 &  0.302 & 0.763\\
            \hline
            Age[60+]:Icon[TO] & -2.156 & -2.696 & 0.007\\
            \hline
        \end{tabular}
        \caption{Results of Speed model 4 with (No websites 1 and 3) or ($>200$s total Speed or $>80$\% average Misclick Rate)}
        \label{fig:placeholder}
    \end{figure}

    The main finding of significance between the 60+ Age Group and the TO Icon Type still holds, with a coefficient of -2.156 and a p-value of 0.007. With the additional exclusion criteria, the time it takes for task completion is still less than any other Icon Type, but the amount that it decreased by is less than in the original linear regression (coefficient=-3.018).

    These findings also mirror the original linear regression in that neither the IWTL interaction term nor either of the Icon Type terms were significant (p-values=0.235, 0.964, 0.763). It is of note that the TO coefficient is positive (coefficient=0.026), albeit not significant, which opposes the results of the original regression (coefficient=-0.264). 

    Now onto Misclick Rate. Like the chosen exclusion criteria where model 3 prevailed, the model of best fit for Misclick Rate is model 3 (8087.388515715943), which barely is lower than model 4's AIC (8089.129). 

    \begin{figure}[H]
        \hypertarget{n13_m3_mr}{}
        \centering
        \begin{tabular}{| p{5cm} | p{3.5cm} | p{3.5cm} | p{3.5cm} |}
            \hline
            & Coefficient & z & P$>|\text{z}|$\\
            \hline
            Icon[IWTL] & 4.158 &  2.133 & 0.033\\
            \hline
            Icon[TO] & -6.044 & -3.262 & 0.001\\
            \hline
        \end{tabular}
        \caption{Results of Misclick Rate model 3 with (No websites 1 and 3) or ($>200$s total Speed or $>80$\% average Misclick Rate)}
        \label{fig:placeholder}
    \end{figure}

    \begin{figure}[H]
        \hypertarget{n13_m4_mr}{}
        \centering
        \begin{tabular}{| p{5cm} | p{3.5cm} | p{3.5cm} | p{3.5cm} |}
            \hline
            & Coefficient & z & P$>|\text{z}|$\\
            \hline
            Icon[IWTL] & 1.951 & 0.726 & 0.468\\
            \hline
            Icon[TO] & -5.591 & -2.123 &  0.034\\
            \hline
            Age[60+]:Icon[IWTL] & 4.611 & 1.189 & 0.234\\
            \hline
            Age[60+]:Icon[TO] & -0.727 & -0.197 & 0.844\\
            \hline
        \end{tabular}
        \caption{Results of Misclick Rate model 4 with (No websites 1 and 3) or ($>200$s total Speed or $>80$\% average Misclick Rate)}
        \label{fig:placeholder}
    \end{figure}

    In model 3, we see that both Icon Type terms hold significance (p-values=0.033, 0.001). This is the same for the linear regression ran in Section \hyperlink{3.3.2}{3.3.2}; where the Misclick Rate model 3 has Icon Type terms with p-values of 0.019 and 0.001. Their coefficients (-6.094, 4.158) behave somewhat as expected-- the original coefficients for IWTL and TO were 6.255 and -8.351, respectively.

    In model 4, IWTL is not significant (p-value=0.468) while TO is (p-value=0.034), which falls in line with the results of model 4 from Section \hyperlink{3.3.2}{3.3.2} (p-values=0.330, 0.037). We also see that both interaction terms do not have significance (p-values=0.234, 0.844), which is the same as the interaction terms from the formal linear regression (p-values=0.331, 0.864). The signs and magnitudes of these coefficients are to expectations. 

\hypertarget{3.3.2}{\subsection{(No websites 1 and 3) or ($>200$s total Speed)}}

    The best-fitting model for Speed was model 4 (AIC=6116.787), which was the same as the linear regression in Section \hyperlink{3.3.2}{3.3.2}.

    \begin{figure}[H]
        \hypertarget{n13_200_speed_m4}{}
        \centering
        \begin{tabular}{| p{5cm} | p{3.5cm} | p{3.5cm} | p{3.5cm} |}
            \hline
            & Coefficient & z & P$>|\text{z}|$\\
            \hline
            Icon[IWTL] & 1.072 & 1.371 & 0.170\\
            \hline
            Icon[TO] & -0.264 & -0.340 & 0.734\\
            \hline
            Age[60+]:Icon[IWTL] & 1.732 &  1.545 & 0.122\\
            \hline
            Age[60+]:Icon[TO] & -1.796 & -1.655 & 0.098\\
            \hline
        \end{tabular}
        \caption{Results of Speed model 4 with (No websites 1 and 3) or ($>200$s total Speed)}
        \label{fig:placeholder}
    \end{figure}

    The main finding no longer holds-- although the p-value for Age[60+]:Icon[TO] is on the lower end at 0.098, it is not significant. The non-significance of the rest of these Icon Type-related terms (p-values=0.170, 0.734, 0.122) are expected given our original findings (see Section \hyperlink{3.2}{3.2}), as well as the signs and magnitudes of the coefficients.

    Given the graph, since only a handful of participants were ($>200$s total Speed or $>80$\% average Misclick Rate), perhaps removing even fewer people out of the mere handful reduced the significance of the interaction. 

    As for Misclick Rate, things look a bit more familiar. Once again, model 3 has the lowest AIC at 8828.632, which is in line with the original linear regression. 

    \begin{figure}[H]
        \hypertarget{n13_200_m3_mr}{}
        \centering
        \begin{tabular}{| p{5cm} | p{3.5cm} | p{3.5cm} | p{3.5cm} |}
            \hline
            & Coefficient & z & P$>|\text{z}|$\\
            \hline
            Icon[IWTL] & 6.910 & 2.603 & 0.009\\
            \hline
            Icon[TO] & -8.059 & -3.137 & 0.002\\
            \hline
        \end{tabular}
        \caption{Results of Misclick Rate model 3 with (No websites 1 and 3) or ($>200$s total Speed)}
        \label{fig:placeholder}
    \end{figure}

    \begin{figure}[H]
        \hypertarget{n13_200_m4_mr}{}
        \centering
        \begin{tabular}{| p{5cm} | p{3.5cm} | p{3.5cm} | p{3.5cm} |}
            \hline
            & Coefficient & z & P$>|\text{z}|$\\
            \hline
            Icon[IWTL] & 3.678 & 0.997 & 0.319\\
            \hline
            Icon[TO] &-7.813 & -2.133 &  0.033\\
            \hline
            Age[60+]:Icon[IWTL] &  6.642 & 1.256  & 0.209\\
            \hline
            Age[60+]:Icon[TO] & -0.252  &  -0.049 &  0.961\\
            \hline
        \end{tabular}
        \caption{Results of Misclick Rate model 4 with (No websites 1 and 3) or ($>200$s total Speed)}
        \label{fig:placeholder}
    \end{figure}

    Once again, we find that in model 3, both IWTL and TO are significant terms (p-values=0.009, 0.002), which matches the findings of the original linear regression. In model 4, significance levels are to expectations as well. Only the TO term is significant (p-value=0.033), whereas IWTL and the rest of the Age Group and Icon Type interaction terms are not (p-values=0.319, 0.209, 0.961). 

    Both model 3 and model 4's terms have coefficients that match the original regression's findings.

\hypertarget{3.3.3}{\subsection{(No websites 1 and 3) or ($>80$\% average Misclick Rate)}}

    Once more, model 4 performs the best for Speed given these criteria, with an AIC of 6330.502. This follows the AICs of Speed in the original regression, where model 4 also had the lowest AIC comparative to the rest of the proposed models (see Section \hyperlink{3.2}{3.2}).

    \begin{figure}[H]
        \hypertarget{n13_80_speed_m4}{}
        \centering
        \begin{tabular}{| p{5cm} | p{3.5cm} | p{3.5cm} | p{3.5cm} |}
            \hline
            & Coefficient & z & P$>|\text{z}|$\\
            \hline
            Icon[IWTL] & 0.787 & 0.842 & 0.400\\
            \hline
            Icon[TO] & -0.127 & -0.139 & 0.890\\
            \hline
            Age[60+]:Icon[IWTL] &  0.830 & 0.634 & 0.526\\
            \hline
            Age[60+]:Icon[TO] & -3.155 & -2.506 & 0.012\\
            \hline
        \end{tabular}
        \caption{Results of Speed Model 4 with (No websites 1 and 3) or ($>80$\% average Misclick Rate)}
        \label{fig:placeholder}
    \end{figure}

    The main findings hold, with a coefficient of -3.155 and a p-value of 0.012. This coefficient is closer to the actual findings, which had a coefficient of -3.018. Like the original findings, none of the other Icon Type-related terms held significance (p-values=0.400, 0.890. 0.526). 

    As for Misclick Rate, we see somewhat familiar trends. This time around, the model of best fit was model 4 (AIC=8874.829), which barely beat out model 3's AIC of 8875.672. This does not follow with the original regression, where model 3 won barely over model 4 (AICs=9048.093, 9050.568). However, the difference between the two AICs for this robustness check's AICs is small, even smaller than the difference in AICs exhibited by the original regression's models for Misclick Rate. As mentioned previously, drops of this magnitude aren't considered to signify a large difference in fitness in models, so this result is unsupported, especially given the findings of both model 3 and 4's findings, which are both provided for completeness.

    \begin{figure}[H]
        \hypertarget{n13_80_m3_mr}{}
        \centering
        \begin{tabular}{| p{5cm} | p{3.5cm} | p{3.5cm} | p{3.5cm} |}
            \hline
            & Coefficient & z & P$>|\text{z}|$\\
            \hline
            Icon[IWTL] & 5.473 & 2.323 & 0.020\\
            \hline
            Icon[TO] & -7.353 & -3.250 & 0.001\\
            \hline
        \end{tabular}
        \caption{Results of Misclick Rate model 3 with (No websites 1 and 3) or ($>80$\% average Misclick Rate)}
        \label{fig:placeholder}
    \end{figure}

    \begin{figure}[H]
        \hypertarget{n13_80_m4_mr}{}
        \centering
        \begin{tabular}{| p{5cm} | p{3.5cm} | p{3.5cm} | p{3.5cm} |}
            \hline
            & Coefficient & z & P$>|\text{z}|$\\
            \hline
            Icon[IWTL] &  1.971 & 0.592 & 0.554\\
            \hline
            Icon[TO] &-5.654 & -1.732 &  0.083\\
            \hline
            Age[60+]:Icon[IWTL] & 6.897 & 1.479 & 0.139\\
            \hline
            Age[60+]:Icon[TO] & -3.046 & -0.679 & 0.497\\
            \hline
        \end{tabular}
        \caption{Results of Misclick Rate model 4 with (No websites 1 and 3) or ($>80$\% average Misclick Rate)}
        \label{fig:placeholder}
    \end{figure}

    Like the original results (see Section \hyperlink{3.2}{3.2}), both IWTL and TO exhibit significance (p-values=0.020, 0.001). The coefficients are within expectations, as well.

    Even though model 4's AIC is lower than model 3's, p-values for these Icon Type-related terms are still over 0.05 (0.554, 0.083, 0.139, 0.497). This mirrors the original findings near-completely, the only exception is that TO's term was found to be significant (p-value=0.037) in the original linear regression. These coefficients also fall in line with previous results.

%     In the following section, the results of the Linear Regression discussed above in \hyperlink{3.2}{3.2} are compared to these regressions performed on similar, alternative exclusion criteria. 

%     After settling upon the exclusion criteria of no Websites 1 and 3, the alternative criteria to confirm robustness were:
%     \begin{enumerate}
%         \item None
%         \item Website 1
%         \item Website 3
%     \end{enumerate}

%     The other options for exclusion criteria mentioned in \hyperlink{2.3.4}{2.3.4} were also tried before running regressions, but deemed not as relevant to the current criteria. These results, however, can still be found in the Appendix.

% \hypertarget{3.3.1}{\subsection{No Exclusions}}
%     Again, following the format of section \hyperlink{3.2}{3.2}, AIC results will first be shown, followed by relevant Linear Regression results.

%     First, we discuss Speed. Model 3's AIC score was the lowest at 11813.533233944247. The model of best fit is the one including all terms except the Age Group and Icon Type interaction term, meaning these interactions terms did not provide helpful information retention for the model. Interestingly enough, AIC actually increased when the variable representing Icon was folded into the equation, meaning it actually did worse than just Age by itself. Concerning, but nevertheless we move towards the results themselves for more conclusive findings. The following are snippets from the results from \hyperlink{ne_m3_speed}{models 3} and \hyperlink{ne_m4_speed}{4} that are most relevant for comparison.

%     \begin{figure}[H]
%         \hypertarget{ne_m3_speed}{}
%         \centering
%         \begin{tabular}{| p{5cm} | p{3.5cm} | p{3.5cm} | p{3.5cm} |}
%             \hline
%             & Coefficient & z & P$>|\text{z}|$\\
%             \hline
%             Icon[IWTL] & 1.612 & 1.343 & 0.179\\
%             \hline
%             Icon[TO] & -0.428 & -0.369 & 0.712\\
%             \hline
%         \end{tabular}
%         \caption{Results of Speed Model 3 with No Exclusions}
%         \label{fig:placeholder}
%     \end{figure}

%     \begin{figure}[H]
%         \hypertarget{ne_m4_speed}{}
%         \centering
%         \begin{tabular}{| p{5cm} | p{3.5cm} | p{3.5cm} | p{3.5cm} |}
%             \hline
%             & Coefficient & z & P$>|\text{z}|$\\
%             \hline
%             Icon[IWTL] & 2.107 & 1.234 & 0.217\\
%             \hline
%             Icon[TO] & 0.527 & 0.310 & 0.756\\
%             \hline
%             Age[60+]:Icon[IWTL] & -0.979 & -0.405 & 0.685\\
%             \hline
%             Age[60+]:Icon[TO] & -1.817 & -0.776 & 0.438\\
%             \hline
%         \end{tabular}
%         \caption{Results of Speed Model 4 with No Exclusions}
%         \label{fig:placeholder}
%     \end{figure}

%     In model 3, Icon Types had no significance. IWTL actually was closer to being significant than TO (p-values=0.179, 0.712), but neither were particularly close to significant levels. As a note, all other terms (Age Group, Websites) exhibited p-values of 0.000.

%     In model 4, neither Icon Types nor the interaction term exhibited any significance. The coefficients for these terms are outliers from what we see throughout the main regression (see \hyperlink{3.2}{3.2}) and the rest of the robustness check results. Both IWTL and TO are positive (coefficients=2.107, 0527), meaning that overall both performed worse than IO. This seems to contradict calculations for \hyperlink{avg_speed}{Average Speed}, but this could be due to model error. However, both interaction terms had negative coefficients (coefficients=-0.979, -1.817), implying better performance from the 60+ group than the 18-35 group in response to these specific Icon Types. This also seems to contradict findings presented in \hyperlink{3.1}{3.1}, but again this could be due to model error.
    
%     Again, IWTL got the closest out of the four terms to significance (p-value=0.217). As another note, all other terms not listed (Age Group, Websites) exhibited p-values of 0.000.

%     As for Misclick Rate, model 3 once again outperforms the rest with an AIC of 13927.39569189007, meaning the interaction terms do not benefit model fitness much. Onto selected results:

%     \begin{figure}[H]
%         \hypertarget{ne_m3_mr}{}
%         \centering
%         \begin{tabular}{| p{5cm} | p{3.5cm} | p{3.5cm} | p{3.5cm} |}
%             \hline
%             & Coefficient & z & P$>|\text{z}|$\\
%             \hline
%             Icon[IWTL] & 4.838 &  1.852 & 0.064\\
%             \hline
%             Icon[TO] & -3.376 & -1.337 & 0.181\\
%             \hline
%         \end{tabular}
%         \caption{Results of Misclick Rate Model 3 with No Exclusions}
%         \label{fig:placeholder}
%     \end{figure}

%     \begin{figure}[H]
%         \hypertarget{ne_m4_mr}{}
%         \centering
%         \begin{tabular}{| p{5cm} | p{3.5cm} | p{3.5cm} | p{3.5cm} |}
%             \hline
%             & Coefficient & z & P$>|\text{z}|$\\
%             \hline
%             Age[60+] & 7.145 & 1.963 & 0.050\\
%             \hline
%             Icon[IWTL] & 2.624 & 0.712 & 0.477\\
%             \hline
%             Icon[TO] & -3.661 & -1.000 & 0.317\\
%             \hline
%             Age[60+]:Icon[IWTL] & 4.447 &  0.853 & 0.394\\
%             \hline
%             Age[60+]:Icon[TO] & 0.639 & 0.127 & 0.899\\
%             \hline
%         \end{tabular}
%         \caption{Results of Misclick Rate Model 4 with No Exclusions}
%         \label{fig:placeholder}
%     \end{figure}

%     In model 3, IWTL eeks towards significance (p-value=0.064), but still doesn't hit the mark. Again, all other coefficients had p-values of 0.000.  
    
%     In model 4, coefficients behave more in line with expectations than when looking at Speed (coefficients=2.624, -3.661, 4.447, 0.639), but still nothing of significance. Interesting to note is that Age is barely no longer significant (p-value=0.050), meaning this model did not notice significantly different error rates between the Age Groups. All website terms had p-values of 0.000.

%     Between both Speed and Misclick Rate on no exclusions, the jarringly inaccurate model reportings showcase how deeply noisy this data was. This was expected, as detailed in \hyperlink{2.1.1}{2.1.1}, but the \hyperlink{ne_m4_speed}{wrongly signed coefficients} and \hyperlink{ne_m4_mr}{nonsignificance of Age} nail this point home.

% \hypertarget{3.3.2}{\subsection{Website 1}}
%     As a reminder, Website 1 was initially excluded alongside Website 3 because it was first in the experiment for all users, and this made it prone to a learning curve. 

%     Once again, the lowest AIC for Speed was model 3 (AIC=9119.597747044225), meaning that the interaction terms did not add significant benefits to model fitness.

%     Models 3 and 4 reaped the following results:

%     \begin{figure}[H]
%         \hypertarget{no_w1_m3_speed}{}
%         \centering
%         \begin{tabular}{| p{5cm} | p{3.5cm} | p{3.5cm} | p{3.5cm} |}
%             \hline
%             & Coefficient & z & P$>|\text{z}|$\\
%             \hline
%             Icon[IWTL] & 1.894 &  2.058 & 0.040\\
%             \hline
%             Icon[TO] & -0.430 & -0.483 & 0.629\\
%             \hline
%         \end{tabular}
%         \caption{Results of Speed Model 3 with no Website 1}
%         \label{fig:placeholder}
%     \end{figure}

%     \begin{figure}[H]
%         \hypertarget{no_w1_m4_speed}{}
%         \centering
%         \begin{tabular}{| p{5cm} | p{3.5cm} | p{3.5cm} | p{3.5cm} |}
%             \hline
%             & Coefficient & z & P$>|\text{z}|$\\
%             \hline
%             Icon[IWTL] & 2.217 & 1.712 &  0.087\\
%             \hline
%             Icon[TO] & 0.883 & 0.686 & 0.493\\
%             \hline
%             Age[60+]:Icon[IWTL] & -0.626 &  -0.342 & 0.733\\
%             \hline
%             Age[60+]:Icon[TO] & -2.482 & -1.398 & 0.162\\
%             \hline
%         \end{tabular}
%         \caption{Results of Speed Model 4 with no Website 1}
%         \label{fig:placeholder}
%     \end{figure}

%     In model 3, the most shocking breakthrough is the significance of IWTL (p-value=0.040). Its coefficient of 1.894 points towards IWTL being worse for speed than IO in general.

%     However, this doesn't hold in model 4. Model 3 has more accuracy in fit, so that explains part of it. But again, the signs on the coefficients are not to expectation. Icon Type terms both have positive signs (coefficients=2.217, 0.883), meaning both perform worse than IO on Websites 2-6. But the interaction terms are both negative (coefficients=-0.626, -2.482), meaning older users perform better given these conditions versus the younger participant pool. But both the Icon Type and interaction terms do not have significance (p-values=0.087, 0.493, 0.733, 0.162). 

%     As for Misclick Rate, once more model 3 outperforms the rest of the batch (AIC=11460.732933369016). So once more, we look towards both models 3 and 4 for insight.

%     \begin{figure}[H]
%         \hypertarget{no_w1_m3_mr}{}
%         \centering
%         \begin{tabular}{| p{5cm} | p{3.5cm} | p{3.5cm} | p{3.5cm} |}
%             \hline
%             & Coefficient & z & P$>|\text{z}|$\\
%             \hline
%             Icon[IWTL] & 7.990 & 3.103 & 0.002\\
%             \hline
%             Icon[TO] & -2.790 & -1.120 & 0.263\\
%             \hline
%         \end{tabular}
%         \caption{Results of Misclick Rate Model 3 with no Website 1}
%         \label{fig:placeholder}
%     \end{figure}

%     \begin{figure}[H]
%         \hypertarget{no_w1_m4_mr}{}
%         \centering
%         \begin{tabular}{| p{5cm} | p{3.5cm} | p{3.5cm} | p{3.5cm} |}
%             \hline
%             & Coefficient & z & P$>|\text{z}|$\\
%             \hline
%             Icon[IWTL] & 5.998 & 1.650 &  0.099\\
%             \hline
%             Icon[TO] & -3.010 & -0.833 & 0.405\\
%             \hline
%             Age[60+]:Icon[IWTL] & 4.000 &   0.778 & 0.437\\
%             \hline
%             Age[60+]:Icon[TO] & 0.507  &  0.102 & 0.919\\
%             \hline
%         \end{tabular}
%         \caption{Results of Misclick Rate Model 4 with no Website 1}
%         \label{fig:placeholder}
%     \end{figure}

%     We see that IWTL remains significant in model 3 (p-value=0.002), matching Speed's similar findings. Once again mirroring Speed, model 4 does not share this finding. In fact, none of the Icon Type and interaction terms hold significance (p-value=0.099, 0.405, 0.437, 0.919). IWTL is the closest out of the four terms to significance (p-value=0.099), though.

%     To recap, IWTL seems to have significance when looking just at Websites 2-6. If the exclusionary criteria reasoning was correct, then Website 1's main issue was being placed first in the user flow. If that's the case, errors could have been distributed in a way that obscured IWTL's speed-inhibiting effects, as it does not rely on the Icon Type encountered if instructions were unclear or zoomed past. 

% \hypertarget{3.3.3}{\subsection{Website 3}}

%     To recap, Website 3 was initially excluded due to the complexity of the task. The UI was crowded by various design elements (see Figure \hyperlink{w3}{2.5}), which is necessary for a food ordering webpage. And the success criteria was clicking a small button amongst the visually busy UI. It's deep enough within the sequence of Websites that it does suffer from the learning curve that Website 1 does, but its generally high difficulty level of a task created a lot of noise in its findings.

%     First, we look at Speed. The lowest AIC belongs to model 3 (AIC=9634.502880155569). This follows a similar trend with the rest of the robustness check's AICs.

%     Now, looking at results from models 3 and 4, we get:

%     \begin{figure}[H]
%         \hypertarget{no_w3_m3_speed}{}
%         \centering
%         \begin{tabular}{| p{5cm} | p{3.5cm} | p{3.5cm} | p{3.5cm} |}
%             \hline
%             & Coefficient & z & P$>|\text{z}|$\\
%             \hline
%             Icon[IWTL] & 1.103 & 1.032 & 0.302\\
%             \hline
%             Icon[TO] & -1.585 & -1.534 & 0.125\\
%             \hline
%         \end{tabular}
%         \caption{Results of Speed Model 3 with no Website 3}
%         \label{fig:placeholder}
%     \end{figure}

%     \begin{figure}[H]
%         \hypertarget{no_w3_m4_speed}{}
%         \centering
%         \begin{tabular}{| p{5cm} | p{3.5cm} | p{3.5cm} | p{3.5cm} |}
%             \hline
%             & Coefficient & z & P$>|\text{z}|$\\
%             \hline
%             Icon[IWTL] & 1.169 & 0.776 & 0.438\\
%             \hline
%             Icon[TO] & -0.463 & -0.309 & 0.757\\
%             \hline
%             Age[60+]:Icon[IWTL] & -0.114 &  -0.053 & 0.958\\
%             \hline
%             Age[60+]:Icon[TO] & -2.112 & -1.023 & 0.306\\
%             \hline
%         \end{tabular}
%         \caption{Results of Speed Model 4 with no Website 3}
%         \label{fig:placeholder}
%     \end{figure}

%     None of the Icon Type terms are significant according to model 3, but it is worth noting that this time, TO's closer to significance than IWTL, as was seen in excluding just Website 1 (see \hyperlink{no_w1_m3_speed}{above}). 

%     Again, none of the Icon Type or interaction terms were significant either. But again, the closest to significance was the interaction term for 60+/TO, which was not the case for excluding just Website 1 (it was the second closest [p-value=0.162] after IWTL [p-value=0.087]). As a sidebar, all other terms in both model 3 and 4 were significant.

%     AIC was no different for Misclick Rate; model 3 still remained the model with the lowest AIC at  11526.779337953953. 

%     As for results, we see that it mirrors a lot of what Speed found:

%     \begin{figure}[H]
%         \hypertarget{no_w3_m3_mr}{}
%         \centering
%         \begin{tabular}{| p{5cm} | p{3.5cm} | p{3.5cm} | p{3.5cm} |}
%             \hline
%             & Coefficient & z & P$>|\text{z}|$\\
%             \hline
%             Icon[IWTL] & 2.820 & 1.013 & 0.311\\
%             \hline
%             Icon[TO] & -7.942 & -2.951 & 0.003\\
%             \hline
%         \end{tabular}
%         \caption{Results of Misclick Rate Model 3 with no Website 3}
%         \label{fig:placeholder}
%     \end{figure}

%     \begin{figure}[H]
%         \hypertarget{no_w3_m4_mr}{}
%         \centering
%         \begin{tabular}{| p{5cm} | p{3.5cm} | p{3.5cm} | p{3.5cm} |}
%             \hline
%             & Coefficient & z & P$>|\text{z}|$\\
%             \hline
%             Icon[IWTL] & 0.092 & 0.024 & 0.981 \\
%             \hline
%             Icon[TO] & -7.634 & -1.959 & 0.050\\
%             \hline
%             Age[60+]:Icon[IWTL] & 5.488 &  0.989 & 0.323\\
%             \hline
%             Age[60+]:Icon[TO] & -0.449 & -0.084 & 0.933\\
%             \hline
%         \end{tabular}
%         \caption{Results of Misclick Rates Model 4 with no Website 3}
%         \label{fig:placeholder}
%     \end{figure}

%     Excitingly, TO holds significance in model 3 (p-value=0.002). IWTL, unlike in the models excluding only Website 1, was not (p-value=0.087). Model 4 also mimics this trend with TO, although TO is barely nonsignificant (p-value=0.005). 

%     Excluding Website 3 brings to light the positive effect of TO, even at times the positive effect it specifically has on the 60+ Age Group. This complements the results of excluding solely Website 1, which only showed the negative effects of IWTL. Perhaps TO was able to shine through because complexity of the task surpassed any one Icon Type's ability to positively influence performance metrics.

%     To summarize, these robustness checks point towards the main results, but do not completely validate them. Without Website 1, the significantly negative impacts of IWTL (p-value for Speed=0.040, p-value for Misclick Rate=0.002) that are mirrored in the main results shine through. Without Website 3, TO's positive impact was signaled towards (p-value for Misclick Rate=0.003). Neither displayed similar significant interactions between the 60+ Age Group and the TO Icon Type. Throughlines are drawn, but the main results do not hold in terms of robustness. Alternatively, these main results are sensitive to specific Websites, namely 2, 4, 5, and 6. This should be kept at front of mind when reading the \hyperlink{4}{Discussion} portion of this thesis. 
    

\hypertarget{4}{\chapter{Discussion}}

    This experiment was intricate and littered with many new challenges for the computer science-trained thesis author. That being said, there is a lot to discuss. 
    
    Generally, the format of this chapter will be as follows: First, in Section \hyperlink{4.1}{4.1}, qualitative conjectures will be attached to quantitative results found above in Chapter \hyperlink{3}{3}. Then, Section \hyperlink{4.2}{4.2} will bring to light, confront, and acknowledge the various limitations of this thesis. Finally, future directions will be suggested in Sections \hyperlink{4.3}{4.3} and \hyperlink{4.4}{4.4}.

\hypertarget{4.1}{\section{Interpretations}}
    Two main results follow from Chapter \hyperlink{3}{3}, namely that the Text-Only condition has a disproportionately positive influence on the 60+ Age Group (and that TO overall performs better across Age Groups than IO and IWTL), and that the Icon with Text Labels condition overall performs worse than TO and IO conditions. A reminder that these results hold for Websites 2, 4, 5, and 6.

\hypertarget{4.1.1}{\subsection{Text-Only and 60+ Interpretations}}
    To reiterate, when modeling Speed, the best model of fit included the Age Group and Icon Type interaction terms. More specifically, the interaction term between the 60+ Age Group and the TO Icon Type had a coefficient of -3.018 and a p-value of 0.018. In other words, the 60+ Age Group averaged 3.018 second less on tasks that were shown in the TO condition than the 18-35 group.

    As for Misclick Rate and TO, in both semi-complete model 3 and fully-expanded model 4, TO by itself had significance (p-values=0.001, 0.037) both with negative coefficients (coefficients=-8.351, -7.813). However, the interaction term from above was not significant in Misclick Rate.

    These findings, especially the significance of the TO and 60+ interaction in relation to Speed, are incredibly exciting. As explained in Section \hyperlink{1}{1}, one of the main focuses of the age-based accessibility subfield of HCI is finding ways to even the technological playing field for older users experiencing age-related physical and mental declines. Of course, finding accessibility fixes that all ages benefit from is also exciting, and in practice can be employed in building accessibility technology. But the grand prize is discovering tweaks that overall benefit both Age Groups but ultimately benefit older users more than younger users. It's in line with notions of equality versus equity-- one makes sure all are given equal resources despite different starting lines, and the other ensures all are given adequate resources to make the playing field completely equal. 

    This result finds backing in scholarship, as well. As mentioned in , the closer the semantic distance icons are to what they represent, the faster participants are. And what closes the semantic distance more than language itself? However, there is debate as to whether semantic distance provides a larger performance boost to older users versus younger (see ). To recap, Liu et al.\@ found that further semantic distance disproportionately hindered the reaction times of older participants, as compared to younger participants \cite{liu_study_2023}; Dosso and Chevalier also found similar results \cite{dosso_how_2021}. But these age-based discrepancies are contested by Ganor and Te'eni \cite{ganor_designing_nodate} and Schröder and Ziefle \cite{schroder_effects_2008}, so this thesis' findings add a new perspective into this inconclusive, longstanding HCI debate.

    But as seen in Chapter \hyperlink{1}{1}, of the literature written on icon accessibility, very few included a Text-Only condition. To some extent, the reason is obvious-- in studying iconography, the goal is to condense large amounts of word-based meaning into a visual iota, usually for UI/UX design purposes. But this suggests exploration into a middle ground between alienating language entirely and alienating design principles completely might hold merit. 

    Of course, this study's results do not come without room for questioning. One of the main conjectures that will be presented to contextualize the negative impact of IWTL below is that the UI experienced high levels of clutter, to use Grahame et al.\@'s terminology \cite{grahame_age_2004}. Although the positive performance effects of TO are safely significant, it could be possible that they benefited from lack of clutter, although this would not account for the significant difference between TO and IO specifically. And another possibility is that \enquote{digital native} 18-35 year olds might not have been used to associating text with buttons, especially without the explicit hyperlink formatting (which the TO condition was not presented in). Without this prior conditioning, it could have been the case that TO would no longer significantly boost older users over younger users. But with significant results in Misclick Rate that extends equally to both groups, this is unlikely. 

\hypertarget{4.1.2}{\subsection{Icon with Text Labels Interpretations}}
    The significant results related to the Icon with Text Labels condition were found within models run with Misclick Rate as the dependent variable. In the model of best fit, model 3, IWTL had a coefficient of 6.225 and a p-value of 0.019. Thus, IWTL made Misclick Rate higher, specifically by 6.225\%, across the board for both Age Groups, as compared to IO and TO. For transparency, these results were not significant in model 4, for either IWTL by itself or the 60+ and IWTL interaction term; IWTL was also not significantly worse in Speed for any model.

    These results directly conflict with the consensus within the literature. Generally, IWTL conditions reap faster and more accurate results as compared to IO conditions. To summarize, Dosso and Chevalier observed that labeled icons gave older participants speed boosts on par with their younger counterparts \cite{dosso_how_2021}. This finding is seconded by Wu et al.\@ \cite{wu_how_2022}. Wu et al.\@ specifically found that skeuomorphic icons with text outperformed its flat icons with text and unlabeled counterparts. Song et al.\@'s work that icons presented with text labels enabled older users to perform at the same accuracy levels as their younger counterparts \cite{song_developing_2024}. There is some contention from Islam, whose combined quantitative and behavioral study concluded there was no frontrunner between the IO and IWTL conditions \cite{islam_exploring_2015}, but the sect within HCI which disagrees with IWTL's positive performance benefits is small.

    So why does IWTL significantly and negatively impact accuracy in this experiment? To propose an explanation, these studies were all conducted in a laboratory setting where icons (or icons with text labels) were shown in isolation. This is not the case in a task-based setting such as this thesis. Icons now exist within a fuller context, which brings along confounding variables such as clutter \cite{grahame_age_2004} that do not exist within a controlled, sterile environment. This is both a combination of a change in environment as well as experimental design error, which especially shone through when combined with already complex and information-dense UIs such as Website 3.

\hypertarget{4.2}{\section{Limitations}}
    Now, we move towards the limitations of this study, some of which were touched upon by Section \hyperlink{4.1}{4.1} already.

    The first order of business is acknowledging faulty experimental design. This existed in terms of UI/UX design choices (see Website 3's results in Section \hyperlink{3.1.1}{3.1.1}), specifically in terms of the general instructions provided to participants. Although much of the error which arose with Website 1 could be attributed with user error, some of it could have originated in unclear instructions. An example of an instruction screen presented to participants is shown in Figure \hyperlink{instructions}{4.1}:

    \begin{figure}[H]
        \hypertarget{instructions}{}
        \centering
        \fbox{
\includegraphics[width=0.75\linewidth]{instructions.png}}
        \caption{An example of an instruction screen provided during the experiment}
    \end{figure}

    This was initially altered to provide more clarity (e.g. with wording such as \enquote{Task 1: Switch to image search} in a larger font), but the original version was shown to participants due to a Maze bug. 

    Another limitation related to experimental design was the decision to host the experiment online and asynchronously through Prolific. There are justifications for this choice (see ) and precedent within both HCI and larger fields that feed into HCI (see ), but this choice also left us with noisy data before exclusions (see Section \hyperlink{3.3}{3.3}) and an astounding drop-off rate of $\sim38.9\%$.

    One of the main questions that arose after seeing the skewed proportion of dropped-off older versus younger users (see , e.g. for IO, 6 18-35 users did not finish while 17 60+ users did not finish) was whether \enquote{tech-savvy} older users were being disproportionately represented in the dataset, as compared to \enquote{non-tech-savvy} older users. Sampling from a pool of already technologically-fluent users was another concern about sampling from Prolific (see ).

    % The reasoning for not including dropped-off user responses within the dataset was that they would have contributed to the already high levels of noise within the dataset. That is to be expected, but when compounded against the noise that was derived from general user error in relation to the online and asynchronous nature of this study, it obscures the signal within heavily. A trend that was qualitatively observed was that users who dropped-off after making it to the UIs would \enquote{give up} after realizing they didn't get the answer immediately. This mirrors a phenomenon outlined in Mahmud et al.\@, where older users exhibit a larger amount of unproductive behaviors stemming from frustration due to obstruction from task-completion as compared to younger Age Groups \cite{mahmud_learning_2020}, which might explain the higher drop-off rate for older users as compared to younger. To some extent, random clicking, erratic mouse movements, etc.\@ imitate real-life digital behaviors, 

    The reasoning for not including dropped-off users' responses within the dataset was that incomplete session did not provide interpretable measures of task performance via Speed or Misclick Rate. Indeed, task abandonment and behaviors indicative of user frustration (such as erratic mouse movements and random clicking) \cite{mahmud_learning_2020} are real-world behaviors of older users and are important to acknowledge. But once participants disengaged from the task, subsequent interaction patterns reflected abandonment behavior rather than purposeful navigation, and most importantly, left us with no measurable DV.
    
    In general, much of this could have been mitigated by conducting an in-person experiment. First and foremost, any percentage of the drop-off that was due to technical error would be eradicated. But with user error, in-person experiments can perform \enquote{quality control} on participants, which ensure their willingness to give a proper attempt, which despite the monetary reward, is not always the norm on online recruitment platforms such as Prolific. User error also could have been decreased with guidance from experiment leads when confusion arose during the tasks as well.

    The last limitation that will be acknowledged is tracking accuracy via Misclick Rate. This is due to its rapid growth at onset and plateau shortly after, which creates poles on both ends and skews reporting statistics such as means and medians. This does not negate its usefulness entirely, especially as a metric which provides insight into an additional dimension of task-completion (accuracy), but it certainly made Section \hyperlink{3}{3} less interpretable. It should be mentioned that this statistic is not without backing within the literature \cite{yu_reducing_2024}, but accuracy in general is not widely reported on as much in part because of an inability to find a more suitable alternative. For instance, converting Misclick Rate into its numerator, the number of incorrect clicks, would still result in two tails, especially with the added phenomenon of random clicking due to frustration that was prevalent in this online study.

\hypertarget{4.3}{\section{Extensions}}
    One of the main lingering questions is whether or not this exciting interaction between TO and 60+ is context-specific. The websites provided in this study were selected mainly for their everyday functionality, and without more specific criteria for UI design (hence why Website 3 had an incredibly complex UI as compared to the others). A future direction that could be pursued is replicating this study with a different set of UIs, or even a more curated sequence of interfaces. 

    Another main extension would be to replicate this study in-person. Achieving the same sample size would be difficult, but would be a worthy endeavor, especially given the limitations of online studies outlined above. This would also fall more in line with the standards for current HCI experimental design, but these will shift as fields that have adopted online studies more completely into their toolkits continue to influence the world of HCI.

    Some other directions that would be less important, but still interesting, to explore would be rotating around the order of appearance of the current UIs, and changing what type of icons are displayed (currently, flat icons are employed [see Section \hyperlink{1.3}{1.3}])-- it could be worth exploring implementing skeuomorphic iconography or alternative forms of iconography instead.

\hypertarget{4.4}{\section{Conclusion}}

    As technology continues to evolve in rapid, perhaps even unsustainable, rates \cite{chong_bad_2023}, the requirements of users need to be centered more than ever before. User-centric design is touted as the golden standard of product development in the technological sphere, but is not an inherently all-inclusive framework. Instead, software development and UI/UX design typically centers the conventional user, someone with fluency in technological navigation and no lack in cognitive, visual, motor, etc.\@ abilities. This myth of the \enquote{everyday user} sidelines many people of varying ability levels, including user populations such as the elderly who experience unique combinations of disabilities associated with aging. 

    But there are reasons technologists return to the \enquote{everyday user} archetype time and time again. Accessibility can be a daunting topic, especially in the context of technology. Users require different modes of assistance, and hitting the goal post for every type of disability with a singular user interface can be an overwhelming task. And the guides that do describe how to comprehensively promote accessibility on your platforms, like the WCAG and ARIA, are incredibly dense and long reads.

    But accessibility can also start small. In fact, icons (or text or icons with text labels, in the case of this thesis) are essentially the smallest unit of meaning on a technological interface. And even these tiny units of meaning can significantly and positively enhance the experience of older users in massive ways (think on the magnitude of seconds, even!). 
    % Age-based accessibility research in human-computer interactions aims to unearth various ways of making the technological experiences of older users easier and more productive. These findings require flexibility, as age-based declines do not have a definitive prototype. 

    % Icon Type is one of the research fields within this discipline which has been hotly contested. Results on whether certain alterations benefit users, such as skeuomorphic versus flat icon design and with or without text labels, are all in flux. It is also a field that has primarily existed within the context of laboratory experimentation. But with the need for more conclusions, especially conclusions that are relevant to the everyday use cases of these icons, there became a gap that this thesis could fill.

    % This thesis explored whether different Icon Types, namely Icon-Only, Text-Only, and Icon with Text Labels, when presented in the context of a web-based task, had different age-based performance effects. Each participant was assigned a Icon Type condition and solved the same sequence of 6 UIs online and asynchronously. 
    
    % A total of 240 users completed this experiment, with roughly ~80 users per Icon Type and ~40 users per Icon Type and Age Group subdivisions. After excluding Websites 1 and 3, the main finding from this study was that the Text-Only condition disproportionately positively impacted the Speed of the 60+ Age Group. Alongside that, Text-Only was proven to be significantly positive for accuracy, and Icon with Text Labels was proven to be significantly negative for accuracy. 

    % Thank you.

    
    

    

    
    

    



















%If you feel it necessary to include an appendix, it goes here.
    \appendix
      \chapter{The First Appendix}
      \chapter{The Second Appendix, for Fun}


%This is where endnotes are supposed to go, if you have them.
%I have no idea how endnotes work with LaTeX.

  \backmatter % backmatter makes the index and bibliography appear properly in the t.o.c...

% if you're using bibtex, the next line forces every entry in the bibtex file to be included
% in your bibliography, regardless of whether or not you've cited it in the thesis.
    \nocite{*}

% Rename my bibliography to be called "Works Cited" and not "References" or ``Bibliography''
% \renewcommand{\bibname}{Works Cited}

%    \bibliographystyle{bsts/mla-good} % there are a variety of styles available; 
%  \bibliographystyle{plainnat}
% replace ``plainnat'' with the style of choice. You can refer to files in the bsts or APA 
% subfolder, e.g. 
 \bibliographystyle{IEEEtran}  % or
 \bibliography{thesis}
 % Comment the above two lines and uncomment the next line to use biblatex-chicago.
 %\printbibliography[heading=bibintoc]

% Finally, an index would go here... but it is also optional.
\end{document}
